\subsection{\texorpdfstring{Compact convergence in \(\mathcal{M}(\mathrm{X};\mathbf{C})\)}{Compact convergence in M(X; C)}}

Recall that the topology of compact convergence on \(\mathcal{M}(\mathrm{X};\mathbf{C})\) is the topology of uniform convergence in the compact subsets of \(\mathcal{K}(\mathrm{X};\mathbf{C})\) . We shall call topology of strictly compact convergence on \(\mathcal{M}(\mathrm{X};\mathbf{C})\) the topology of uniform convergence in the strictly compact subsets (No.~1) of \(\mathcal{K}(\mathrm{X};\mathbf{C})\) .

PROPOSITION 17. --- On the space \(\mathcal{M}(\mathrm{X};\mathbf{C})\) , consider the following topologies:

\(\mathcal{T}_1\) : the topology of pointwise convergence in a total subset \(\mathrm{T}\) of \(\mathcal{K}(\mathrm{X};\mathbf{C})\) ;

\(\mathcal{T}_2\) : the vague topology;

\(\mathcal{T}_3\) : the topology of strictly compact convergence;

\(\mathcal{T}_4\) : the topology of compact convergence.

Each of these topologies is coarser than the next. Moreover:

\begin{enumerate}
\def\labelenumi{(\roman{enumi})}
\item
  The bounded sets are the same for \(\mathcal{T}_2\) , \(\mathcal{T}_3\) and \(\mathcal{T}_4\) .
\item
  If \(\mathbf{H}\) is a vaguely bounded subset of \(\mathcal{M}(\mathrm{X};\mathbf{C})\) , the topologies induced on \(\mathbf{H}\) by the topologies \(\mathcal{T}_1, \mathcal{T}_2, \mathcal{T}_3, \mathcal{T}_4\) are identical.
\end{enumerate}

A vaguely bounded subset H of \(\mathcal{M}(X;C)\) is equicontinuous (No.~9, Prop. 15), thus the first assertion follows from TVS, III, §3, No.~7, Prop. 9, and the second from GT, X, §2, No.~4, Th. 1.

Recall that when X is paracompact, the topology of strictly compact convergence coincides with the topology of compact convergence (No.~1, Prop. 2).

PROPOSITION 18. --- On the cone \(\mathcal{M}_{+}(X)\) , the topologies induced by the following topologies coincide:

\(\mathcal{T}_{1}\) : the topology of pointwise convergence in a linear subspace V of \(\mathcal{K}(X;C)\) that is dense in \(\mathcal{K}(X;C)\) and satisfies the property (P) (No.~7, Prop. 9);

\(\mathcal{T}_2\) : the vague topology;

\(\mathcal{T}_3\) : the topology of strictly compact convergence.

Since every filter is the intersection of the ultrafilters finer than it (GT, I, §6, No.~4, Prop. 7), it suffices to show that if \(\mathfrak{U}\) is an ultrafilter on \(\mathcal{M}_{+}(\mathrm{X})\) that converges to a measure \(\mu_0\) for the topology \(\mathcal{T}_1\) , then it also converges to \(\mu_0\) for \(\mathcal{T}_3\) . Let \(\mathbf{K}\) be a compact subset of \(\mathbf{X}\) ; by hypothesis, there exists a function \(h \in \mathrm{V}\) that is \(\geqslant 0\) on \(\mathbf{X}\) and takes values \(>0\) on \(\mathbf{K}\) ; it follows that every function \(f \in \mathcal{K}(\mathrm{X},\mathrm{K};\mathbf{C})\) may be written \(f = gh\) with \(g \in \mathcal{K}(\mathrm{X},\mathrm{K};\mathbf{C})\) , and if \(c = \inf_{x \in \mathrm{K}} h(x) > 0\) then \(\| g \| \leqslant c^{-1} \| f \|\) . By hypothesis, there exists a set \(\mathrm{H}_0 \in \mathfrak{U}\) such that, for every measure \(\mu \in \mathrm{H}_0\) ,

\[
0 \leqslant \mu (h) \leqslant \mu_ {0} (h) + 1 = b.
\]

Consequently, for every function \(f \in \mathcal{K}(\mathrm{X};\mathbf{C})\) ,

\[
| \langle f, h \cdot \mu \rangle | = | \langle h f, \mu \rangle | \leqslant \| f \| \cdot \mu (h) \leqslant b \| f \|
\]

for every measure \(\mu\in\mathrm{H}_{0}\) ; this proves that the set H of measures \(h\cdot\mu\) , where \(\mu\) runs over \(\mathrm{H}_{0}\) , is vaguely bounded. If \(\mathfrak{U}_{0}\) is the ultrafilter induced by \(\mathfrak{U}\) on \(\mathrm{H}_{0}\) , the image of \(\mathfrak{U}_{0}\) under the mapping \(\mu\mapsto h\cdot\mu\) is the base of an ultrafilter \(\mathfrak{F}\) on H, and since H is relatively compact for the topology of strictly compact convergence (Prop. 17 and No.~9, Prop. 15), \(\mathfrak{F}\) is convergent to a measure \(\nu_{0}\) for this topology. In other words, for any \(\varepsilon>0\) and any compact subset L of \(\mathcal{K}(\mathrm{X},\mathrm{K};\mathbf{C})\) , there exists a subset N of \(\mathrm{H}_{0}\) belonging to \(\mathfrak{U}\) such that, for every function \(g\in\mathrm{L}\) and every pair of measures \(\mu,\mu'\) belonging to N, one has \(|\langle g,h\cdot\mu\rangle-\langle g,h\cdot\mu'\rangle|\leqslant\varepsilon\) , that is,

\[
| \langle g h, \mu \rangle - \langle g h, \mu^ {\prime} \rangle | \leqslant \varepsilon .
\]

Now, we saw above that the mapping \(g \mapsto gh\) is an automorphism of the Banach space \(\mathcal{K}(\mathrm{X},\mathrm{K};\mathbf{C})\) . We have thus shown that U is a Cauchy filter on \(\mathcal{M}_{+}(\mathrm{X})\) for the topology of strictly compact convergence. A fortiori, it is a Cauchy filter for vague convergence, and Prop. 14 of No.~9 shows that it is vaguely convergent to a measure \(\mu_{1}\) ; moreover, since V is dense in \(\mathcal{K}(\mathrm{X};\mathbf{C})\) , the hypothesis implies that \(\mu_{1} = \mu_{0}\) ; finally, since U is a Cauchy filter for the topology of strictly compact convergence, it also converges to \(\mu_{0}\) for this topology (GT, X, §1, No.~5, Prop. 5).

Q.E.D.

COROLLARY. --- If X is paracompact then the topologies induced on \(\mathcal{M}_{+}(X)\) by the vague topology and the topology of compact convergence coincide.

However, the topologies induced on \(\mathcal{M}_{+}(X)\) by the topology of compact convergence and the topology of strictly compact convergence may be different when X is not paracompact (Exer. 3).

\section{SUPPORT OF A MEASURE}

\subsection{Restriction of a measure to an open set. Definition of a measure by means of local data}

Let X be a locally compact space, Y an open set in X. The subspace Y of X is locally compact, and every continuous function with values in a topological vector space E, defined on Y and with compact support, may be extended by continuity to all of X, by giving it the value 0 on \(\complement \mathrm{Y}\); one can therefore identify in this way the space \(\mathcal{K}(\mathrm{Y};\mathrm{E})\) with the linear subspace of \(\mathcal{K}(\mathrm{X};\mathrm{E})\) formed by the continuous functions with compact support contained in Y. If \(\mu\) is a measure on X, it is clear that the restriction of \(\mu\) to \(\mathcal{K}(\mathrm{Y};\mathbf{C})\) is a measure on Y, which is called the restriction of \(\mu\) to the open subspace Y, or the measure induced on Y by \(\mu\) , and is denoted \(\mu|\mathrm{Y}\) . The restrictions to Y of \(|\mu|\) , \(\mathcal{R}\mu\) and \(\mathcal{I}\mu\) are, respectively, \(|\mu|\mathrm{Y}|\) , \(\mathcal{R}(\mu|\mathrm{Y})\) and \(\mathcal{I}(\mu|\mathrm{Y})\) , by virtue of §1, Nos. 5 and 6. If \(\mu\) is real then the restrictions of \(\mu^{+}\) and \(\mu^{-}\) to Y are, respectively, \((\mu|\mathrm{Y})^{+}\) and \((\mu|\mathrm{Y})^{-}\) , by virtue of formula (8) of §1, No.~5.

One sees immediately that if Y and Z are two open sets in X such that \(Y \supset Z\) , and if \(\mu|Y\) and \(\mu|Z\) are the restrictions of \(\mu\) to Y and Z, then \(\mu|Z\) is also the restriction of \(\mu|Y\) to the open subspace Z of the locally compact space Y.

In Ch. IV, §5, No.~7 we shall generalize this definition to the case that Y is a locally compact subspace of X.

\pandocbounded{\includegraphics[keepaspectratio]{images/part001_8e0e9629.jpg}}

Note that a measure on Y is not necessarily the restriction of some measure on X (cf.~Ch. V, §7, No.~2, Prop. 3).

For example, let Y be the open interval {]}0,1{[} of X = R; the mapping

\[
f \mapsto \int_ {0} ^ {1} \frac {f (x)}{x} d x
\]

is a measure on Y, because every function in \(\mathcal{K}(\mathrm{Y};\mathbf{C})\) is zero on a neighborhood of 0 in \(\mathbf{R}\) . However, this measure cannot be extended to a measure on \(\mathbf{R}\) because, in the contrary case, its restriction to the set of functions \(f\in \mathcal{K}(\mathrm{Y};\mathbf{C})\) such that \(\| f\| \leqslant 1\) would be bounded; but this is false.

However, we have the following proposition:

PROPOSITION 1. --- Let \((\mathrm{Y}_{\alpha})_{\alpha\in\mathrm{A}}\) be an open covering of X and suppose given, on each subspace \(Y_{\alpha}\) , a measure \(\mu_{\alpha}\) , in such a way that for every pair \((\alpha,\beta)\) , the restrictions of \(\mu_{\alpha}\) and \(\mu_{\beta}\) to \(Y_{\alpha}\cap Y_{\beta}\) are identical. Under these conditions, there exists one and only one measure \(\mu\) on X whose restriction to \(Y_{\alpha}\) is equal to \(\mu_{\alpha}\) for every index \(\alpha\) .

Let us first show that every function \(f \in \mathcal{K}(\mathrm{X};\mathbf{C})\) may be written in the form of a finite sum \(f = \sum_{i} f_{i}\) where, for each of the functions \(f_{i} \in \mathcal{K}(\mathrm{X};\mathbf{C})\) , there exists an index \(\alpha_{i}\) such that \(\operatorname{Supp}(f_{i}) \subset \mathrm{Y}_{\alpha_{i}}\) . If \(\mathrm{K} = \operatorname{Supp}(f)\) , there exists a finite number of indices \(\alpha_{i} (1 \leqslant i \leqslant n)\) such that the \(\mathrm{Y}_{\alpha_{i}}\) form a covering of \(\mathrm{K}\) ; let \(h_{i} (1 \leqslant i \leqslant n)\) be continuous mappings of \(\mathrm{X}\) into {[}0,1{]} such that the support of \(h_{i}\) is compact and is contained in \(\mathrm{Y}_{\alpha_{i}}\) for \(1 \leqslant i \leqslant n\) , and such that \(\sum_{i=1}^{n} h_{i}(x) = 1\) on \(\mathrm{K} (\S 1, \text{No. 2, Lemma 1})\) ; the functions \(f_{i} = f h_{i}\) meet the requirements. This shows in the first place that if there exists a measure \(\mu\) meeting the requirements then it is unique, because for every finite sum \(f = \sum_{i=1}^{n} f_{i}\) , where \(f_{i} \in \mathcal{K}(\mathrm{Y}_{\alpha_{i}}; \mathbf{C})\) , necessarily \(\mu(f) = \sum_{i=1}^{n} \mu_{\alpha_{i}}(f_{i})\) . Moreover, we will have shown the existence of a linear form \(\mu\) on \(\mathcal{K}(\mathrm{X};\mathbf{C})\) whose restriction to each subspace \(\mathcal{K}(\mathrm{Y}_{\alpha};\mathbf{C})\) is \(\mu_{\alpha}\) , provided we demonstrate the following property: if \((g_{i})_{1 \leqslant i \leqslant m}\) and \((h_{j})_{1 \leqslant j \leqslant n}\) are two finite sequences of functions in \(\mathcal{K}(\mathrm{X};\mathbf{C})\) such that \(g_{i} \in \mathcal{K}(\mathrm{Y}_{\alpha_{i}}; \mathbf{C})\) for \(1 \leqslant i \leqslant m\) , \(h_{j} \in \mathcal{K}(\mathrm{Y}_{\beta_{j}}; \mathbf{C})\) for \(1 \leqslant j \leqslant n\) and

\[
\sum_ {i = 1} ^ {m} g _ {i} (x) = \sum_ {j = 1} ^ {n} h _ {j} (x) = 1
\]

on K, then

\[
\sum_ {i = 1} ^ {m} \mu_ {\alpha_ {i}} (f g _ {i}) = \sum_ {j = 1} ^ {n} \mu_ {\beta_ {j}} (f h _ {j}).
\]

Now,

\[
f g _ {i} = \sum_ {j = 1} ^ {n} f g _ {i} h _ {j},
\]

whence

\[
\sum_ {i = 1} ^ {m} \mu_ {\alpha_ {i}} (f g _ {i}) = \sum_ {i = 1} ^ {m} \sum_ {j = 1} ^ {n} \mu_ {\alpha_ {i}} (f g _ {i} h _ {j}).
\]

Similarly,

\[
\sum_ {j = 1} ^ {n} \mu_ {\beta_ {j}} (f h _ {j}) = \sum_ {j = 1} ^ {n} \sum_ {i = 1} ^ {m} \mu_ {\beta_ {j}} (f g _ {i} h _ {j}).
\]

But since the support of \(fg_{i}h_{j}\) is contained in \(\mathbf{Y}_{\alpha_i}\cap \mathbf{Y}_{\beta_j}\) , we have \(\mu_{\alpha_i}(fg_i h_j) = \mu_{\beta_j}(fg_i h_j)\) , which establishes our assertion.

It remains to see that \(\mu\) is a measure on X; now, every point of X admits a compact neighborhood contained in one of the \(Y_{\alpha}\) ; the conclusion therefore follows at once from the definition of \(\mu\) and from Prop. 6 of §1, No.~3.

COROLLARY (Principle of localization). --- Let \(\mu\) and \(\nu\) be two measures on X, and let \((\mathrm{Y}_{\alpha})\) be a family of open sets of X such that, for every \(\alpha\) , the restrictions to \(Y_{\alpha}\) of \(\mu\) and \(\nu\) are equal; then the restrictions of \(\mu\) and \(\nu\) to \(Y = \bigcup_{\alpha} Y_{\alpha}\) are equal.

\subsection{Support of a measure}

Let \(\mu\) be a measure on a locally compact space \(X\) , and let \(\mathfrak{G}\) be the set of open sets \(U \subset X\) such that the restriction of \(\mu\) to \(U\) is zero; it follows at once from the principle of localization (No.~1, Cor. of Prop. 1) that if \(U_0\) is the union of the sets \(U \in \mathfrak{G}\) , then \(U_0\) itself belongs to \(\mathfrak{G}\) and is therefore the largest of the sets of \(\mathfrak{G}\) .

DEFINITION 1. --- If \(\mu\) is a measure on a locally compact space \(X\) , one defines the support of \(\mu\) , denoted \(\operatorname{Supp}(\mu)\) , to be the closed set complementary to the largest of the open sets in \(X\) on which the restriction of \(\mu\) is zero.

To say that a point \(x \in X\) does not belong to the support of \(\mu\) means that there exists an open neighborhood V of x such that the restriction of \(\mu\) to V is zero; to say that x belongs to the support of \(\mu\) therefore means that for every neighborhood V of x, there exists a function \(f \in \mathcal{K}(X; \mathbf{C})\) , whose support is contained in V, such that \(\mu(f) \neq 0\) .

Examples. --- 1) For a measure on X to be zero, it is necessary and sufficient that its support be empty.

\begin{enumerate}
\def\labelenumi{\arabic{enumi})}
\setcounter{enumi}{1}
\item
  The support of Lebesgue measure on \(\mathbf{R}\) is the entire line \(\mathbf{R}\) ; for, it is nonempty and is invariant under every translation.
\item
  On the interval \(\mathrm{X} = [0,1]\) of \(\mathbf{R}\) consider a countable dense subset, arranged into a sequence \((a_n)\) , and let \(\mu\) be the measure defined by placing the mass \(2^{-n}\) at the point \(a_n\) for every \(n \geqslant 0\) (§1, No.~3, Example I). The support of \(\mu\) is all of \(\mathrm{X}\) ; for, let \(x\) be any point of \(\mathrm{X}\) , \(\mathrm{V}\) a neighborhood of \(x\) , and \(f\) a continuous real-valued function \(\geqslant 0\) on \(\mathrm{X}\) , equal to 1 at the point \(x\) , whose support is contained in \(\mathrm{V}\) (§1, No.~2, Lemma 1); the set of \(y \in \mathrm{V}\) such that \(f(y) > 0\) is open in \(\mathrm{X}\) , therefore contains a point \(a_n\) , consequently \(\mu(f) \geqslant f(a_n)2^{-n} > 0\) .
\end{enumerate}

PROPOSITION 2. --- The support of a measure \(\mu\) is identical to the support of the measure \(|\mu|\) ; if \(\mu\) is real, its support is the union of the supports of the measures \(\mu^{+}\) and \(\mu^{-}\) .

For, if the restriction of \(\mu\) to an open set U is zero, then so is that of \(|\mu|\) (resp. of \(\mu^{+}\) and \(\mu^{-}\) when \(\mu\) is real), and conversely.

\pandocbounded{\includegraphics[keepaspectratio]{images/part001_dfbe8ad7.jpg}}

Note that the supports of \(\mu^{+}\) and \(\mu^{-}\) can be nonempty and identical (cf.~Ch. V, § 5, Exer. 5).

PROPOSITION 3. --- If \(\mu\) and \(\nu\) are two measures on a locally compact space X such that \(|\mu| \leqslant |\nu|\) , then \(\operatorname{Supp}(\mu) \subset \operatorname{Supp}(\nu)\) .

For, if the restriction of \(\nu\) to an open set is zero, then so is that of \(\mu\) .

PROPOSITION 4. --- The support of the sum of two measures is contained in the union of their supports.

For, if the restrictions of two measures to an open set are zero, then the same is true of the restriction of their sum.

If \(\mu\) and \(\nu\) are two positive measures, then the support of \(\lambda = \mu + \nu\) is equal to the union of the supports of \(\mu\) and \(\nu\) ; for, if \(x_0\) is a point of the union and V is any neighborhood of \(x_0\) , then there exists a continuous function \(f \geqslant 0\) with support contained in V and such that one of the numbers \(\mu(f)\) , \(\nu(f)\) is \(>0\) ; a fortiori, \(\lambda(f) = \mu(f) + \nu(f) > 0\) .

PROPOSITION 5. --- The support of the restriction of a measure \(\mu\) to an open set U is the trace on U of the support of \(\mu\) .

The proposition is obvious from the definitions.

PROPOSITION 6. --- The set of measures on a locally compact space X, whose support is contained in a closed set F, is a vaguely closed linear subspace of \(\mathcal{M}(\mathrm{X};\mathbf{C})\) .

For, it is the intersection of the vaguely closed hyperplanes with equation \(\mu(f) = 0\) , where f runs over the set of functions in \(\mathcal{K}(\mathrm{X};\mathbf{C})\) whose support does not intersect F.

Suppose X is not compact: given a filter \(\Phi\) on the space \(\mathcal{M}(\mathrm{X};\mathbf{C})\) of measures on X, we shall say that the support of a measure \(\mu\) recedes indefinitely along \(\Phi\) if, for every compact subset K of X, there exists a set \(M\in\Phi\) such that \(\operatorname{Supp}(\mu)\cap K=\varnothing\) for every measure \(\mu\in M\) .

PROPOSITION 7. --- If \(\Phi\) is a filter on \(\mathcal{M}(\mathrm{X};\mathbf{C})\) such that the support of \(\mu\) recedes indefinitely along \(\Phi\) , then \(\mu\) converges vaguely to 0 with respect to \(\Phi\) .

Let \(f\) be any function in \(\mathcal{K}(\mathrm{X};\mathbf{C})\) and let \(K\) be its support. By hypothesis, there exists a set \(M \in \Phi\) such that for every measure \(\mu \in M\) , \(\operatorname{Supp}(\mu) \cap K = \varnothing\) ; it follows that \(\mu(f) = 0\) for all \(\mu \in M\) , which proves the proposition.

\subsection{Characterization of the support of a measure}

By definition, if the support of a function \(f \in \mathcal{K}(\mathrm{X};\mathbf{C})\) does not intersect the support of a measure \(\mu\) , then \(\mu(f) = 0\) ; but the following more precise result is true:

PROPOSITION 8. --- Let \(\mu\) be a measure on a locally compact space \(X\) . For every function \(f \in \mathcal{K}(X; \mathbf{C})\) that is zero on \(\operatorname{Supp}(\mu)\) , \(\mu(f) = 0\) .

Set \(\mathrm{K} = \operatorname{Supp}(f)\) , \(\mathrm{S} = \operatorname{Supp}(\mu)\) . Given a number \(\varepsilon > 0\) , let \(\mathrm{V}\) be the set of \(x \in \mathrm{X}\) such that \(|f(x)| < \varepsilon\) ; \(\mathrm{V}\) is an open set containing \(\mathrm{S}\) by hypothesis, therefore \(\complement \mathrm{S}\) is a neighborhood of the compact set \(\complement \mathrm{V}\) . It follows that there exists a continuous mapping \(h\) of \(\mathrm{X}\) into {[}0,1{]}, equal to 1 on \(\complement \mathrm{V}\) and with support contained in \(\complement \mathrm{S}\) (§1, No.~2, Lemma 1). Since the support of \(fh\) does not intersect \(\mathrm{S}\) , \(\mu(fh) = 0\) . On the other hand, \(f = fh\) on \(\mathrm{K} \cap \complement \mathrm{V}\) , and \(|fh| \leqslant |f|\) on \(\mathrm{X}\) , therefore \(|f - fh| \leqslant 2\varepsilon\) on \(\mathrm{X}\) , by the choice of \(\mathrm{V}\) . Observe finally that there exists a number \(\mathrm{M}_{\mathrm{K}}\) such that \(|\mu(g)| \leqslant \mathrm{M}_{\mathrm{K}} \| g\|\) for every function \(g \in \mathcal{K}(\mathrm{X};\mathbf{C})\) whose support is contained in \(\mathrm{K}\) ; since the support of \(f - fh\) is contained in \(\mathrm{K}\) , \(|\mu(f - fh)| \leqslant 2\mathrm{M}_{\mathrm{K}}\varepsilon\) and consequently \(|\mu(f)| = |\mu(f - fh)| \leqslant 2\mathrm{M}_{\mathrm{K}}\varepsilon\) ; since \(\varepsilon\) is arbitrary, \(\mu(f) = 0\) .

COROLLARY 1. --- If two functions \(f, g\) in \(\mathcal{K}(\mathrm{X};\mathbf{C})\) are equal on \(\operatorname{Supp}(\mu)\) , then \(\mu(f) = \mu(g)\) .

COROLLARY 2. --- Let \(\mu\) be a positive measure on X; if \(f \in \mathcal{K}(X; \mathbf{C})\) is such that \(f(x) \geqslant 0\) on \(\operatorname{Supp}(\mu)\) , then \(\mu(f) \geqslant 0\) .

For, \(f = |f|\) on \(\operatorname{Supp}(\mu)\) , therefore \(\mu(f) = \mu(|f|) \geqslant 0\) by Corollary 1.

COROLLARY 3. --- Let \(\mu\) be a bounded measure on X; if \(f \in \mathcal{K}(X; \mathbf{C})\) is such that \(|f(x)| \leqslant a\) on \(\operatorname{Supp}(\mu)\) , then \(|\mu(f)| \leqslant a \| \mu \|\) .

For, \(\operatorname{Supp}(|\mu|) = \operatorname{Supp}(\mu)\) , and if \(h\) is a continuous mapping of X into \([0,1]\) , equal to 1 on \(\operatorname{Supp}(f)\) and with compact support, then \(|f(x)| \leqslant ah(x)\) on \(\operatorname{Supp}(\mu)\) , therefore

\[
| \mu | (| f |) \leqslant a | \mu | (h) \leqslant a \| \mu \|
\]

by Corollary 2; the conclusion then follows from formula (13) of §1, No.~6.

PROPOSITION 9. --- Let \(\mu\) be a positive measure on X; if \(f\) is a function in \(\mathcal{K}_{+}(X)\) such that \(\mu(f)=0\) , then \(f\) is zero on \(\operatorname{Supp}(\mu)\) .

Let \(x\) be a point of \(X\) such that \(f(x) > 0\) ; let us show that \(x\) does not belong to \(\operatorname{Supp}(\mu)\) . For, there exist a compact neighborhood \(V\) of \(x\) and a number \(a > 0\) such that \(f(y) \geqslant a\) on \(V\) . If \(g\) is any continuous function \(\geqslant 0\) with support contained in \(V\) , let us show that \(\mu(g) = 0\) ; indeed, if one sets \(b = \|g\|\) then \(g \leqslant bf/a\) , whence \(\mu(g) \leqslant b\mu(f)/a = 0\) .

PROPOSITION 10. --- Let \(\mu\) be a measure on a locally compact space \(X\) ; for every function \(g \in \mathcal{C}(X; \mathbf{C})\) , the support of the measure \(g \cdot \mu\) is the closure \(T\) of the set of points \(x \in \operatorname{Supp}(\mu)\) such that \(g(x) \neq 0\) .

Set \(S = \operatorname{Supp}(\mu)\) ; let \(x_{0}\) be a point not belonging to T; there exists an open neighborhood V of \(x_{0}\) such that at every point of \(V \cap S\) , g is zero; if \(f \in \mathcal{K}(X; \mathbf{C})\) has support contained in V, then fg is zero on S, therefore (Prop. 8) \(\mu(gf) = 0\) ; in other words, the restriction of \(g \cdot \mu\) to V is zero.

Conversely, assuming that the restriction of \(g \cdot \mu\) to an open neighborhood W of a point \(x_{0} \in X\) is zero, let us show that there does not exist a point of \(W \cap S\) at which g is \(\neq 0\) . Indeed, if there were such a point y, there would exist a compact neighborhood U of y, contained in W, at every point x of which \(g(x) \neq 0\) ; but then every function \(f \in \mathcal{K}(X; \mathbf{C})\) with support contained in U could be written f = gh, where \(h \in \mathcal{K}(X; \mathbf{C})\) has support contained in \(U \subset W\) ; it would then follow that \(\mu(f) = \mu(gh) = 0\) , contrary to the hypothesis \(y \in S\) .

\pandocbounded{\includegraphics[keepaspectratio]{images/part001_a1e78629.jpg}}

Note that T is contained in the intersection of the support S of \(\mu\) and the support of g, but it is not necessarily equal to this intersection. For example, if X = R, \(\mu\) is the Dirac measure at the point 0, and \(g(x) = x\) , then \(g \cdot \mu = 0\) even though the intersection of the supports of g and \(\mu\) reduces to the point 0, thus is nonempty.

COROLLARY. --- In order that the measure \(g \cdot \mu\) be zero, it is necessary and sufficient that g be zero on the support of \(\mu\) .

PROPOSITION 11. --- Every measure with compact support is bounded.

For, \(|\mu|\) is also a measure with compact support, thus we can restrict attention to the case that \(\mu \geqslant 0\) ; if \(h\) is a continuous mapping of \(X\) into {[}0,1{]}, with compact support and equal to 1 on \(\operatorname{Supp}(\mu)\) , then for every function \(f \in \mathcal{K}(X; \mathbf{C})\) one has \(|f(x)| \leqslant \|f\|h(x)\) on \(\operatorname{Supp}(\mu)\) , therefore (Cor. 2 of Prop. 8) \(\mu(|f|) \leqslant \mu(h)\|f\|\) , which proves the proposition (§1, No.~8).

\subsection{Point measures. Measures with finite support}

PROPOSITION 12. --- Let \(a_i\) ( \(1 \leqslant i \leqslant n\) ) be distinct points in a locally compact space X. Every measure on X whose support is contained in the set of the \(a_i\) is a linear combination of the measures \(\varepsilon_{a_i}\) ( \(1 \leqslant i \leqslant n\) ).

For, such a measure \(\mu\) is zero for every function \(f\in\mathcal{K}(\mathrm{X};\mathbf{C})\) satisfying the \(n\) relations \(f(a_{i})=0\) (No.~3, Prop. 8); since these relations may be written \(\varepsilon_{a_{i}}(f)=0\) , \(\mu\) is a linear combination of the \(\varepsilon_{a_{i}}\) (A, II, §7, No.~5, Cor. 1 of Th. 7).

In particular, every measure whose support is either empty or reduced to a single point x is of the form \(\alpha\varepsilon_{x}\) , where \(\alpha\) is a complex number; such a measure is said to be a point measure; thus every measure whose support is finite is a sum of point measures.

THEOREM 1. --- Every measure \(\mu\) on a locally compact space \(X\) is in the vague closure of the vector space \(V\) of measures whose support is finite and contained in \(\operatorname{Supp}(\mu)\) .

It suffices to prove that \(\mu\) is orthogonal to the subspace \(V^{\circ}\) of \(\mathcal{K}(X; \mathbf{C})\) orthogonal to \(V\) (TVS, II, §6, No.~3, Cor. 2 of Th. 1), that is, that the relations \(\langle f, \varepsilon_{a} \rangle = 0\) , where \(a\) runs over the support of \(\mu\) , imply \(\langle f, \mu \rangle = 0\) ; but this is just Prop. 8 of No.~3.

COROLLARY 1. --- Every bounded measure \(\mu\) on X is in the vague closure of the convex set A of measures whose support is finite and is contained in that of \(\mu\) , and whose norm is \(\leqslant \| \mu \|\) . Moreover, if \(\nu\) tends vaguely to \(\mu\) while remaining in A, then \(\| \nu \|\) tends to \(\| \mu \|\) .

To prove the first assertion, it suffices to establish that the measure \(\mu\) belongs to the polar set of the polar set \(A^{\circ}\) of A in \(\mathcal{K}(X;C)\) (TVS, II, §6, No.~3, Th. 1 and §8, No.~4); this means that for \(f \in \mathcal{K}(X;C)\) , the relations \(|\langle f,\varepsilon_{a}\rangle| \leqslant 1/\|\mu\|\) for all \(a \in \operatorname{Supp}(\mu)\) imply that \(|\langle f,\mu\rangle| \leqslant 1\) ; but this is a consequence of Cor. 3 of Prop. 8 of No.~3.

To prove the second assertion, we note that

\[
\liminf _ {\nu \to \mu ,   \nu \in A} \| \nu \| \geqslant \| \mu \|
\]

since the function \(\nu \mapsto \| \nu \|\) is lower semi-continuous for the vague topology (§1, No.~9, Cor. 4 of Prop. 15), and the conclusion follows from the fact that \(\| \nu \| \leqslant \| \mu \|\) for \(\nu \in A\) by definition.

COROLLARY 2. --- Every bounded measure \(\mu\) on \(X\) is in the vague closure of the set of measures whose support is finite and contained in that of \(\mu\) and whose norm is equal to \(\|\mu\|\) .

We can suppose that \(\mu \neq 0\) . Let \(V\) be an open neighborhood of 0 for the vague topology; for every \(\varepsilon\) such that \(0 < \varepsilon < 1\) , there exists, by virtue of Cor. 1, a measure \(\nu_0\) whose support is finite and contained in \(\operatorname{Supp}(\mu)\) and for which \(\nu_0 - \mu \in V\) and \(\| \mu \| \geqslant \| \nu_0 \| \geqslant (1 - \varepsilon) \| \mu \|\) . Setting \(\nu = (\| \mu \| / \| \nu_0 \|) \nu_0\) , we have \(\| \nu \| = \| \mu \|\) and \(\| \nu - \nu_0 \| \leqslant \| \mu \|\) ; for \(\varepsilon\) sufficiently small we therefore have \(\nu - \mu \in V + V\) , whence the conclusion.

COROLLARY 3. --- Every bounded positive measure \(\mu\) on \(X\) is in the vague closure of the convex set of positive measures whose support is finite and contained in that of \(\mu\) and whose norm is equal to \(\|\mu\|\) .

The same reasoning as in Cor. 2 shows that we can limit ourselves to proving that \(\mu\) is in the vague closure of the convex set B formed by the positive measures with finite support contained in \(\mathrm{Supp}(\mu)\) and with norm \(\leqslant \| \mu \|\) . Again, it suffices to establish that \(\mu\) belongs to the polar set of \(\mathbf{B}^{\circ}\) , the polar set of B in the space \(\mathcal{K}(\mathbf{X};\mathbf{R})\) (TVS, II, §6, No.~3, Th. 1); but this means that for \(f\in \mathcal{K}(\mathbf{X};\mathbf{R})\) the relations \(\langle f,\varepsilon_a\rangle \geqslant -1 / \| \mu \|\) for all \(a\in \mathrm{Supp}(\mu)\) imply \(\langle f,\mu \rangle \geqslant -1\) , which is a consequence of No.~3, Cor. 2 of Prop. 8.

COROLLARY 4. --- In the space \(\mathcal{M}(\mathrm{X};\mathbf{C})\) , the set of point measures is total for the topology of strictly compact convergence (§1, No.~10).

On the cone \(\mathcal{M}_{+}(X)\) , the topology of strictly compact convergence is identical to the vague topology (§1, No.~10, Prop. 18), and every measure on X may be written \(\mu_{1}-\mu_{2}+i\mu_{3}-i\mu_{4}\) , where the \(\mu_{j}\ (1\leqslant j\leqslant4)\) are positive measures; the conclusion therefore follows from Th. 1.

PROPOSITION 13. --- Let \(\mu\) be a measure on a locally compact space \(X\) . For a point \(x_0\) to belong to \(\operatorname{Supp}(\mu)\) , it is necessary and sufficient that the point measure \(\varepsilon_{x_0}\) be in the vague closure of the set of measures \(g \cdot \mu\) , where \(g\) runs over the set of continuous functions with compact support such that \(\| g \cdot \mu \| \leqslant 1\) .

The condition is obviously sufficient by virtue of Prop. 6 of No.~2. To see that it is necessary, suppose \(x_0 \in \operatorname{Supp}(\mu)\) ; consider a finite number of functions \(f_k\) ( \(1 \leqslant k \leqslant n\) ) in \(\mathcal{K}(X; \mathbf{C})\) , and an arbitrary number \(\delta > 0\) ; we are to prove that there exists a function \(g \in \mathcal{K}(X; \mathbf{C})\) such that \(\| g \cdot \mu \| \leqslant 1\) and such that

\[
\left| f _ {k} \left(x _ {0}\right) - \mu \left(g f _ {k}\right) \right| \leqslant \delta
\]

for \(1 \leqslant k \leqslant n\) . Let U be a relatively compact open neighborhood of \(x_0\) such that the oscillation of each of the \(f_k\) ( \(1 \leqslant k \leqslant n\) ) on U is \(\leqslant \delta/2\) . By hypothesis, since \(x_0 \in \text{Supp}(\mu)\) , there exists a function \(g_0 \in \mathcal{K}(X; \mathbf{C})\) whose support is contained in U and such that \(\mu(g_0) \neq 0\) ; the measure \(\nu = g_0 \cdot \mu\) is not zero, since for every function \(f \in \mathcal{K}(X; \mathbf{C})\) equal to 1 on U, \(\nu(f) = \mu(g_0) \neq 0\) . Moreover, \(\nu\) is bounded (No.~3, Prop. 11);

multiplying \(g_{0}\) by a scalar, we can suppose that \(\|\nu\|=1\) . This being so, setting \(\alpha_{k}=f_{k}(x_{0})\) we can write, for \(1\leqslant k\leqslant n\) and for every function \(h\in\mathcal{K}(\mathrm{X};\mathbf{C})\) ,

\[
f _ {k} (x _ {0}) - \nu (f _ {k} h) = \alpha_ {k} (1 - \nu (h)) + \nu ((\alpha_ {k} - f _ {k}) h).
\]

Since \(\nu\) has its support in U, we may identify it with its restriction to U; the hypothesis \(\| \nu \| = 1\) then implies that there exists a function \(h \in \mathcal{K}(\mathrm{X};\mathbf{C})\) , with support contained in U, such that \(\| h \| \leqslant 1\) and such that \(|\alpha_{k}(1 - \nu (h))| \leqslant \delta / 2\) for \(1 \leqslant k \leqslant n\) . The definition of U moreover shows that \(|(\alpha_{k} - f_{k}(x))h(x)| \leqslant \delta / 2\) for all \(x \in \mathrm{U}\) ; since \(\| \nu \| = 1\) and \(\operatorname{Supp}(\nu) \subset \mathrm{U}\) we therefore have \(|\nu ((\alpha_{k} - f_{k})h)| \leqslant \delta / 2\) and so, setting \(g = g_0h\) ,

\[
\left| f _ {k} \left(x _ {0}\right) - \mu \left(g f _ {k}\right) \right| \leqslant \delta \quad \text { for } 1 \leqslant k \leqslant n.
\]

This proves the proposition, since \(\| g\cdot \mu \| = \| (g_0h)\cdot \mu \| \leqslant \| g_0\cdot \mu \| = 1\)

COROLLARY. --- Let \(\mu\) be a measure on X. For a measure \(\nu\) on X to be in the vague closure of the set of measures \(g \cdot \mu\) , where \(g\) runs over \(\mathcal{K}(X; \mathbf{C})\) , it is necessary and sufficient that \(\operatorname{Supp}(\nu) \subset \operatorname{Supp}(\mu)\) .

For, \(\operatorname{Supp}(g\cdot\mu)\subset\operatorname{Supp}(\mu)\) by No.~3, Prop. 10; therefore the support of every vague limit of measures of the form \(g\cdot\mu\) is also contained in \(\operatorname{Supp}(\mu)\) (No.~2, Prop. 6). Conversely, if \(\operatorname{Supp}(\nu)\subset\operatorname{Supp}(\mu)\) then \(\nu\) is the vague limit of measures with finite support contained in \(\operatorname{Supp}(\mu)\) (Th. 1), hence is in the vague closure of the set of measures \(g\cdot\mu\) by Prop. 13.

\subsection{Discrete measures}

PROPOSITION 14. --- For a measure \(\mu\) on a locally compact space \(X\) to be a discrete measure (§1, No.~3, Example I), it is necessary and sufficient that \(\operatorname{Supp}(\mu)\) be a discrete closed subspace of \(X\) .

Let \(\mu\) be a discrete measure on X, defined by the masses \(h(x) \neq 0\) placed at the points \(x\) of a discrete closed subspace N of X; let us show that \(\operatorname{Supp}(\mu) = \mathbf{N}\) . For every \(a \in \mathbf{N}\) and every neighborhood V of \(a\) , there exists a function \(f \in \mathcal{K}(\mathbf{X}; \mathbf{C})\) with support contained in V, equal to 1 at the point \(a\) and to 0 at the other points of N, whence \(\mu(f) = h(a) \neq 0\) . On the other hand if \(b \notin \mathbf{N}\) then there exists a neighborhood W of \(b\) not intersecting N; for every function \(g \in \mathcal{K}(\mathbf{X}; \mathbf{C})\) with support contained in W, we therefore have \(\mu(g) = 0\) , which proves that \(b \notin \operatorname{Supp}(\mu)\) .

Conversely, let \(\mu\) be a measure such that \(\operatorname{Supp}(\mu)\) is a discrete closed subspace N of X. For every \(a \in N\) , there exists an open neighborhood \(V_{a}\) of a that contains no point of N distinct from a; the restriction of \(\mu\) to \(V_{a}\) is therefore a point measure with support \(\{a\}\) (No.~2, Prop. 5), hence (No.~4, Prop. 12) is of the form \(h(a)\varepsilon_{a}\) , where \(h(a)\neq0\) . Setting \(h(x)=0\) at the points of \(\complement N\), and denoting by \(\nu\) the measure defined by the masses \(h(x)\) , the principle of localization shows that \(\nu=\mu\) .

We thus see that on a discrete space X, every measure is discrete.

\section{INTEGRALS OF CONTINUOUS VECTOR-VALUED FUNCTIONS}

Throughout this section, X denotes a locally compact space, E a locally convex space over R or C. We denote by \(E'\) the dual of E (the space of continuous linear forms on E) and by \(E'^{*}\) the algebraic dual of \(E'\) (the space of all linear forms on \(E'\) ); for \(z \in E\) , \(z' \in E'\) , \(z'^{*} \in E'^{*}\) , we write \(\langle z, z' \rangle = z'(z)\) , \(\langle z'^{*}, z' \rangle = z'^{*}(z')\) .

Recall that if E is Hausdorff, then E may be identified with a linear subspace of \(E^{\prime*}\) by identifying an element \(z \in E\) with the linear form \(z' \mapsto \langle z, z' \rangle\) on \(E'\) , and that \(E^{\prime*}\) , equipped with the weak topology \(\sigma(E^{\prime*}, E')\) , may be canonically identified with the completion of E equipped with the weakened topology \(\sigma(E, E')\) .

\subsection{Definition of the integral of a vector-valued function}

Recall that a mapping \(\mathbf{f}\) of X into E is said to be weakly continuous if, for every \(\mathbf{z}' \in \mathrm{E}'\) , the mapping \(x \mapsto \langle \mathbf{f}(x), \mathbf{z}' \rangle\) of X into \(\mathbf{C}\) (in other words the mapping \(\mathbf{z}' \circ \mathbf{f}\) , also denoted \(\langle \mathbf{f}, \mathbf{z}' \rangle\) ) is continuous. We shall say that a mapping \(\mathbf{f}\) of X into E is scalarly of compact support if, for every \(\mathbf{z}' \in \mathrm{E}'\) , the mapping \(x \mapsto \langle \mathbf{f}(x), \mathbf{z}' \rangle\) has compact support. We denote by \(\widetilde{\mathcal{K}}(\mathrm{X};\mathrm{E})\) the space of mappings of X into E that are weakly continuous and scalarly of compact support; it is clear that \(\widetilde{\mathcal{K}}(\mathrm{X};\mathrm{E}) \supset \mathcal{K}(\mathrm{X};\mathrm{E})\) , but these two spaces are not necessarily identical (see Example 2 below); they are, however, equal when E is finite-dimensional.

Note that in the definition of a function that is weakly continuous (resp. scalarly of compact support), the topology of E intervenes only through the intermediation of the dual \(E'\) of E; thus, the set of these functions is not changed when the topology of E is replaced by any locally convex topology for which the dual is the same.

If E and F are two vector spaces in duality, we note that it means the same to say that a mapping of X into E is continuous for \(\sigma(\mathrm{E},\mathrm{F})\) and to say that it is weakly continuous.

Let \(\mathbf{f}\) be a mapping of \(X\) into \(E\) , weakly continuous and scalarly of compact support, and let \(\mu\) be a measure on \(X\) ; for every \(\mathbf{z}' \in E'\) we have \(\mathbf{z}' \circ \mathbf{f} \in \mathcal{K}(X)\) ; set

\[
\varphi (\mathbf {z} ^ {\prime}) = \int \langle \mathbf {f} (x), \mathbf {z} ^ {\prime} \rangle d \mu (x) = \mu (\mathbf {z} ^ {\prime} \circ \mathbf {f}).
\]

It is clear that \(\varphi\) is a linear form on \(\mathbf{E}'\) , hence is an element of \(\mathbf{E}^{\prime*}\) .

DEFINITION 1. --- For every function \(\mathbf{f} \in \widetilde{\mathcal{K}}(\mathrm{X};\mathrm{E})\) we call integral of \(\mathbf{f}\) with respect to \(\mu\) , and denote by \(\int \mathbf{f} d\mu\) or \(\int \mathbf{f}(x) d\mu(x)\) , or \(\int \mathbf{f}\mu\) , or \(\int \mathbf{f}(x)\mu(x)\) , the element of \(\mathrm{E}'^*\) defined by

\[
\left\langle \int \mathbf {f} d \mu , \mathbf {z} ^ {\prime} \right\rangle = \int \left\langle \mathbf {f}, \mathbf {z} ^ {\prime} \right\rangle d \mu \quad \text {for all} \mathbf {z} ^ {\prime} \in \mathrm{E} ^ {\prime}.\tag{1}
\]

We note that even if E is Hausdorff and \(\mathbf{f} \in \mathcal{K}(\mathrm{X};\mathrm{E})\) , one does not necessarily have \(\int f d\mu \in E\) (Exer. 1; cf.~No.~3).

Examples. --- 1) Suppose that E is finite-dimensional over C and Hausdorff, so that if \((\mathbf{e}_{i})_{1\leqslant i\leqslant n}\) is a basis of E then the mapping

\[
\left(\xi_ {1}, \dots \xi_ {n}\right) \mapsto \sum_ {i = 1} ^ {n} \xi_ {i} \mathbf {e} _ {i}
\]

is an isomorphism of \(\mathbf{C}^n\) onto \(\mathbf{E}\) . We then know that every linear form on \(\mathbf{E}\) is continuous, in other words that \(\mathbf{E}'\) is identical to the algebraic dual \(\mathbf{E}^*\) of \(\mathbf{E}\) , and \(\mathbf{E}'^*\) may be canonically identified with \(\mathbf{E}\) . Let \((\mathbf{e}_i')_{1 \leqslant i \leqslant n}\) be the basis of \(\mathbf{E}'\) dual to \((\mathbf{e}_i)\) ; for a mapping \(\mathbf{f}\) of \(\mathbf{X}\) into \(\mathbf{E}\) to be weakly continuous and scalarly of compact support, it is necessary and sufficient that the functions \(f_i = \mathbf{e}_i' \circ \mathbf{f}\) be continuous with compact support; we then have \(\mathbf{f}(x) = \sum_{i=1}^{n} f_i(x) \mathbf{e}_i\) for all \(x \in \mathbf{X}\) , and

\[
\int \mathbf {f} d \mu = \sum_ {i = 1} ^ {n} \mu (f _ {i}) \mathbf {e} _ {i}.
\]

\begin{enumerate}
\def\labelenumi{\arabic{enumi})}
\setcounter{enumi}{1}
\tightlist
\item
  Let us take for E the space \(\mathcal{M}(\mathrm{X};\mathbf{C})\) of measures on X, equipped with the vague topology ( \(\S1\) , No.~9); the dual \(E'\) of E may then be canonically identified with the space \(\mathcal{K}(\mathrm{X};\mathbf{C})\) (TVS, II, \(\S6\) , No.~2, Prop. 3). The mapping \(x \mapsto \varepsilon_{x}\) of X into E is continuous ( \(\S1\) , No.~9, Prop. 13), but its support is not compact if X is not compact; however, it is scalarly of compact support, because for every function \(f \in E'\) the function \(x \mapsto \langle \varepsilon_x, f \rangle = f(x)\) by definition has compact support. Moreover,
\end{enumerate}

\[
\int \left\langle \varepsilon_ {x}, f \right\rangle d \mu = \int f (x) d \mu (x) = \left\langle \mu , f \right\rangle
\]

for every function \(f \in \mathcal{K}(\mathrm{X};\mathbf{C}) = \mathrm{E}'\) , which proves that

\[
\int \varepsilon_ {x} d \mu (x) = \mu
\]

for every measure \(\mu\) on X.

\begin{enumerate}
\def\labelenumi{\arabic{enumi})}
\setcounter{enumi}{2}
\tightlist
\item
  If \(\mathbf{E}\) is Hausdorff then, for every point \(y \in \mathbf{X}\) and every function \(\mathbf{f} \in \widetilde{\mathcal{K}}(\mathbf{X};\mathbf{E})\) , we have
\end{enumerate}

\[
\int \mathbf {f} d \varepsilon_ {y} = \mathbf {f} (y)
\]

because \(\int \langle \mathbf{f},\mathbf{z}^{\prime}\rangle d\varepsilon_{y} = \langle \mathbf{f}(y),\mathbf{z}^{\prime}\rangle\) by definition.

Remarks. --- 1) If E is a locally convex space and N is the closure of \(\{0\}\) in E, so that \(E_{1} = E/N\) is the Hausdorff locally convex space associated with E, we know that the duals \(E'\) and \(E_{1}'\) are identical; for a function f to belong to \(\widetilde{\mathcal{K}}(X;E)\) , it is necessary and sufficient that \(f_{1} = \pi \circ f\) (where \(\pi : E \to E_{1}\) is the canonical homomorphism) belong to \(\widetilde{\mathcal{K}}(X;E_{1})\) , in which case \(\int f d\mu = \int f_{1} d\mu\) . We may therefore limit ourselves to considering only Hausdorff locally convex spaces.

\begin{enumerate}
\def\labelenumi{\arabic{enumi})}
\setcounter{enumi}{1}
\tightlist
\item
  Let E be a locally convex space over C, and let \(E_{0}\) be the locally convex space over R underlying E; we know that the mapping \(z^{\prime} \mapsto Rz^{\prime}\) which, to every continuous (complex) linear form \(z^{\prime}\) on E, makes correspond the continuous (real) linear form \(z \mapsto R\langle z, z^{\prime} \rangle\) on \(E_{0}\) , is an R-isomorphism of the dual \(E^{\prime}\) onto the dual \(E_{0}^{\prime}\) of \(E_{0}\) (TVS, II, §8, No.~1). Similarly, the algebraic dual \(E_{0}^{\prime *}\) of the real vector space \(E_{0}^{\prime}\) may be canonically identified with the real space underlying the algebraic dual \(E^{\prime *}\) of \(E^{\prime}\) . It follows that if \(\mu\) is a real measure and f a mapping in \(\widetilde{\mathcal{K}}(X;E)\) , the formula (1) is again valid when f is regarded as taking its values in \(E_{0}\) and the canonical bilinear forms figuring in the two members as being, respectively, relative to the duality between \(E_{0}^{\prime}\) and \(E_{0}^{\prime *}\) for the first member and the duality between \(E_{0}\) and \(E_{0}^{\prime}\) for the second.
\end{enumerate}

\subsection{Properties of the vectorial integral}

PROPOSITION 1. --- The mapping

\[
(\mathbf {f}, \mu) \mapsto \int \mathbf {f} d \mu
\]

of \(\widetilde{\mathcal{K}} (\mathrm{X};\mathrm{E})\times \mathcal{M}(\mathrm{X};\mathbf{C})\) into \(\mathrm{E}^{\prime *}\) is bilinear.

The proposition follows immediately from Def. 1 of No.~1.

Let \(\mathbf{u}\) be a continuous linear mapping of \(\mathrm{E}\) into a locally convex space \(\mathrm{F}\) ; we know that the transpose \(^t\mathbf{u}\) is a linear mapping of the dual \(\mathrm{F}'\) of \(\mathrm{F}\) into the dual \(\mathrm{E}'\) of \(\mathrm{E}\) ; we shall denote by \(^{tt}\mathbf{u}\) the mapping \(\mathrm{E}^{\prime *} \to \mathrm{F}^{\prime *}\) , the transpose of \(^t\mathbf{u}\) (in the algebraic sense); when \(\mathrm{E}\) and \(\mathrm{F}\) are Hausdorff and are canonically identified with subspaces of \(\mathrm{E}^{\prime *}\) and \(\mathrm{F}^{\prime *}\) , respectively, \(^{tt}\mathbf{u}\) extends the mapping \(\mathbf{u}\) . With these notations:

PROPOSITION 2. --- Let u be a continuous linear mapping of E into a locally convex space F; for every function \(\mathbf{f} \in \widetilde{\mathcal{K}}(\mathrm{X};\mathrm{E})\) , the function \(\mathbf{u} \circ \mathbf{f}\) belongs to \(\widetilde{\mathcal{K}}(\mathrm{X};\mathrm{F})\) and

\[
\int \mathbf {u} (\mathbf {f} (x)) d \mu (x) = ^ {t t} \mathbf {u} \left(\int \mathbf {f} (x) d \mu (x)\right).\tag{2}
\]

For every continuous linear form \(z^{\prime} \in F^{\prime}\) , we have \(z^{\prime} \circ u \circ f = y^{\prime} \circ f\) on setting \(y^{\prime} = z^{\prime} \circ u = {}^{t}u(z^{\prime}) \in E^{\prime}\) , whence the first assertion; moreover, for all \(z^{\prime} \in F^{\prime}\) ,

\[
\begin{array}{r l} \left\langle \int (\mathbf {u} \circ \mathbf {f}) d \mu , \mathbf {z} ^ {\prime} \right\rangle & = \int \langle \mathbf {u} \circ \mathbf {f}, \mathbf {z} ^ {\prime} \rangle d \mu = \int \langle \mathbf {f}, ^ {t} \mathbf {u} (\mathbf {z} ^ {\prime}) \rangle d \mu \\ & = \left\langle \int \mathbf {f} d \mu , ^ {t} \mathbf {u} (\mathbf {z} ^ {\prime}) \right\rangle = \left\langle {} ^ {t t} \mathbf {u} \left(\int \mathbf {f} d \mu\right), \mathbf {z} ^ {\prime} \right\rangle , \end{array}
\]

whence the formula (2).

PROPOSITION 3. --- For every function \(g \in \mathcal{C}(\mathrm{X};\mathbf{C})\) and every function \(\mathbf{f} \in \widetilde{\mathcal{K}}(\mathrm{X};\mathrm{E})\) , the function \(g\mathbf{f}\) belongs to \(\widetilde{\mathcal{K}}(\mathrm{X};\mathrm{E})\) and

\[
\int \mathbf {f} d (g \cdot \mu) = \int \mathbf {f} g d \mu .\tag{3}
\]

Indeed, for every \(z^{\prime} \in E^{\prime}\) , \(\langle gf, z^{\prime} \rangle = g\langle f, z^{\prime} \rangle\) , whence the first assertion; moreover,

\[
\begin{array}{r l} \left\langle \int \mathbf {f} d (g \cdot \mu), \mathbf {z} ^ {\prime} \right\rangle & = \int \langle \mathbf {f}, \mathbf {z} ^ {\prime} \rangle d (g \cdot \mu) = \int g \langle \mathbf {f}, \mathbf {z} ^ {\prime} \rangle d \mu \\ & = \int \langle g \mathbf {f}, \mathbf {z} ^ {\prime} \rangle d \mu = \left\langle \int g \mathbf {f} d \mu , \mathbf {z} ^ {\prime} \right\rangle , \end{array}
\]

whence (3).

PROPOSITION 4. --- Let \(\mu\) be a positive measure on X, S its support, and f a function in \(\widetilde{\mathcal{K}}(\mathrm{X};\mathrm{E})\) . Suppose E is Hausdorff, and equip the space \(\mathrm{E}^{\prime*}\) with the weak topology \(\sigma(\mathrm{E}^{\prime*},\mathrm{E}^{\prime})\) .

\begin{enumerate}
\def\labelenumi{(\roman{enumi})}
\item
  The integral \(\int \mathbf{f}d\mu\) belongs to the closure C in \(\mathrm{E}^{\prime *}\) of the convex cone generated by \(\mathbf{f}(\mathrm{S})\) .
\item
  If \(\mu\) is bounded, then the integral \(\int \mathbf{f} d\mu\) belongs to the set \(\| \mu \| \cdot D\) , where \(D\) is the closed convex envelope of \(\mathbf{f}(S)\) in \(\mathrm{E}'^*\) .
\end{enumerate}

If \(\mathrm{E}\) is complex, we equip \(\mathrm{E}\) with its underlying real vector space structure, which, as we have seen, does not modify the formula (1).

\begin{enumerate}
\def\labelenumi{(\roman{enumi})}
\tightlist
\item
  We know that C is the intersection of the closed half-spaces in \(E^{\prime*}\) that contain \(\mathbf{f}(\mathrm{S})\) and are determined by closed hyperplanes passing through 0; it therefore suffices to prove that, for \(z^{\prime} \in E^{\prime}\) , the relation \(\langle\mathbf{f}(x), \mathbf{z}^{\prime}\rangle \geqslant 0\) for all \(x \in S\) implies
\end{enumerate}

\[
\left\langle \int \mathbf {f} d \mu , \mathbf {z} ^ {\prime} \right\rangle \geqslant 0;
\]

but since

\[
\left\langle \int \mathbf {f} d \mu , \mathbf {z} ^ {\prime} \right\rangle = \int \left\langle \mathbf {f}, \mathbf {z} ^ {\prime} \right\rangle d \mu ,
\]

this follows from §2, No.~3, Cor. 2 of Prop. 8.

\begin{enumerate}
\def\labelenumi{(\roman{enumi})}
\setcounter{enumi}{1}
\tightlist
\item
  We know that D is the intersection of the closed half-spaces in \(E^{\prime*}\) that contain \(\mathbf{f}(S)\) ; it therefore suffices to show that, for \(z^{\prime} \in E^{\prime}\) , the relation \(\langle\mathbf{f}(x), \mathbf{z}^{\prime}\rangle \leqslant a\) for all \(x \in S\) implies
\end{enumerate}

\[
\left\langle \int \mathbf {f} d \mu , \mathbf {z} ^ {\prime} \right\rangle \leqslant a \| \mu \|;
\]

but this follows from §2, No.~3, Cor. 3 of Prop. 8.

COROLLARY. --- Let \(\mu\) be a positive measure on X, f a mapping belonging to \(\mathcal{K}(\mathrm{X};\mathrm{E})\) , and D the closed convex envelope of \(\mathbf{f}(\mathrm{X})\) in \(\mathrm{E}'^*\) . There exists a number \(a > 0\) such that \(\int \mathbf{f} d\mu \in a \cdot \mathrm{D}\) .

If \(\nu\) is any measure on \(X\) , there exist numbers \(a_1, a_2, a_3, a_4 > 0\) such that \(\int \mathbf{f} d\nu \in a_1\mathrm{D} - a_2\mathrm{D} + ia_3\mathrm{D} - ia_4\mathrm{D}\) .

Suppose first that \(\mu\) is positive; by hypothesis, the support K of f is compact; if \(\nu\) is the restriction of \(\mu\) to a relatively compact open neighborhood of K, then \(\nu\) is bounded and \(\int f d\mu = \int f d\nu \in \|\nu\| \cdot D\) by Prop. 4, (ii). The second result follows from this, since any complex measure may be written as \(\mu_{1} - \mu_{2} + i\mu_{3} - i\mu_{4}\) , where the \(\mu_{j}\) are positive.

PROPOSITION 5. --- Suppose that the space X is compact, and let f be a continuous mapping of X into a Hausdorff locally convex space E. The closed convex envelope of \(\mathbf{f}(\mathbf{X})\) in \(E^{\prime*}\) (for \(\sigma(E^{\prime*}, E^{\prime})\) ) is equal to the set of vectors \(\int f d\mu\) for all of the positive measures \(\mu\) on X of total mass 1.

Let C be the closed convex envelope of \(\mathbf{f}(\mathbf{X})\) in \(E^{\prime*}\) ; since \(\mathbf{f}(\mathbf{X})\) is compact and \(E^{\prime*}\) is complete, C is compact. We already know (Prop. 4)

that \(\int \mathbf{f}d\mu \in \mathbb{C}\) for every measure \(\mu\) belonging to the convex set H of positive measures on X of total mass equal to 1. On the other hand, H is convex and compact for the vague topology (§1, No.~9, Cor. 3 of Prop. 15) and is the closure (for this topology) of the convex set \(\mathrm{H}_0\) of positive measures of mass 1 and finite support (§2, No.~4, Cor. 3 of Th. 1). Now, the image of \(\mathrm{H}_0\) under the mapping \(\mu \mapsto \int \mathbf{f}d\mu\) is the convex envelope \(\mathrm{C}_0\) of \(\mathbf{f}(\mathbf{X})\) in \(\mathrm{E}'^*\) . On the other hand, this mapping is continuous for the vague topology on \(\mathcal{M}(\mathrm{X};\mathbf{C})\) and the topology \(\sigma (\mathrm{E}^{\prime *},\mathrm{E}^{\prime})\) on \(\mathrm{E}'^*\) since \(\langle \int \mathbf{f}d\mu ,\mathbf{z}'\rangle = \int \langle \mathbf{f},\mathbf{z}'\rangle d\mu\) by definition; thus the image of \(\mathrm{H} = \overline{\mathrm{H}}_0\) is a compact convex set containing \(\mathrm{C}_0\) and contained in C; since \(\mathrm{C} = \overline{\mathrm{C}}_0\) , this image is equal to C.

PROPOSITION 6. --- For every continuous mapping with compact support f of X into a Hausdorff locally convex space E, every continuous semi-norm q on E and every measure \(\mu\) on X such that \(\int f d\mu \in E\) ,

\[
q \left(\int \mathbf {f} d \mu\right) \leqslant \int (q \circ \mathbf {f}) d | \mu |.\tag{4}
\]

Let D be the set of \(z \in E\) such that \(q(\mathbf{z}) \leqslant 1\) ; D is closed, convex and contains 0, therefore \(D = D^{\circ\circ}\) (TVS, II, §6, No.~3, Cor. 3 of Th. 1). It therefore suffices to prove that for every \(z' \in D^{\circ}\) ,

\[
\left| \left\langle \int \mathbf {f} d \mu , \mathbf {z} ^ {\prime} \right\rangle \right| \leqslant \int (q \circ \mathbf {f}) d | \mu |;
\]

since

\[
\left\langle \int \mathbf {f} d \mu , \mathbf {z} ^ {\prime} \right\rangle = \int \left\langle \mathbf {f}, \mathbf {z} ^ {\prime} \right\rangle d \mu ,
\]

and since, by the definition of \(D^{\circ}\) ,

\[
| \langle \mathbf {f} (x), \mathbf {z} ^ {\prime} \rangle | \leqslant q (\mathbf {f} (x))
\]

for all \(x \in \mathbf{X}\) , the desired inequality follows from the inequality (13) of §1, No.~6.

\subsection{Criteria for the integral to belong to E}

PROPOSITION 7. --- Let E be a Hausdorff locally convex space, and \(\mathbf{f} \in \mathcal{K}(\mathrm{X};\mathrm{E})\) . If \(\mathbf{f}(\mathrm{X})\) is contained in a complete convex subset A of E, then \(\int \mathbf{f} d\mu \in \mathbf{E}\) .

Let K be the support of f, which is compact by hypothesis. Since f is zero on X - K, \(\mathbf{f}(X)\) is equal to \(\mathbf{f}(K)\) or to \(\mathbf{f}(K) \cup \{0\}\) , therefore is compact since f is continuous and E is Hausdorff; the closed convex envelope C of \(\mathbf{f}(X)\) in E is then precompact (for the uniform structure induced by that of E) (TVS, II, §4, No.~1, Prop. 3). But since C is a closed subset of the complete space A, C is complete and therefore compact; a fortiori, C is compact for the weakened topology \(\sigma(\mathrm{E}, \mathrm{E}')\) ; but since that topology is induced by \(\sigma(\mathrm{E}'*, \mathrm{E}')\) , C is the closed convex envelope of \(\mathbf{f}(X)\) in \(E'\) for the topology \(\sigma(\mathrm{E}'*, \mathrm{E}')\) ; the proof is therefore concluded by the Corollary of Prop. 4 of No.~2.

COROLLARY 1. --- Let E be a Hausdorff locally convex space; for every function \(\mathbf{f} \in \mathcal{K}(\mathrm{X};\mathrm{E})\) , \(\int \mathbf{f} d\mu\) belongs to the completion \(\widehat{\mathbf{E}}\) of E.

Since the duals of E and \(\widehat{E}\) are identical, it suffices to apply Prop. 7 while regarding f as taking its values in \(\widehat{E}\) .

COROLLARY 2. --- If E is a quasi-complete Hausdorff locally convex space, then \(\int \mathbf{f} d\mu \in \mathrm{E}\) for every function \(\mathbf{f} \in \mathcal{K}(\mathrm{X};\mathrm{E})\) .

As noted at the beginning of the proof of Prop. 7, \(\mathbf{f}(\mathrm{X})\) is compact and its closed convex envelope C in E is precompact, hence bounded; but since the set C is closed and bounded, it is complete by hypothesis, and it suffices to apply Prop. 7.

We shall see, in Ch. VI, §1, No.~2, other criteria for \(\int \mathbf{f} d\mu\) to belong to E, which apply in particular to the functions in \(\widetilde{\mathcal{K}}(\mathrm{X};\mathrm{E})\) and not just those in \(\mathcal{K}(\mathrm{X};\mathrm{E})\) .

Corollary 2 of Proposition 7 may be applied in the following two cases: \(1^{\circ}\) E is a Banach space; \(2^{\circ}\) E is the dual of a barreled Hausdorff locally convex space G, and E is equipped with an S-topology, where S is a covering of G by bounded subsets (TVS, III, §4, No.~2, Cor. 4 of Th. 1). For example, Cor. 2 of Prop. 7 can be applied when E is the weak dual of a Banach space, or a space of measures \(\mathcal{M}(\mathrm{Y};\mathbf{C})\) equipped with the vague topology.

If \(\mathbf{X} = \mathbf{R}\) , \(\mu\) is Lebesgue measure on \(\mathbf{R}\) , and \(\mathrm{E}\) is a Banach space, then the integral \(\int \mathbf{f} d\mu\) of a function in \(\mathcal{K}(\mathrm{X};\mathrm{E})\) is none other than the integral

\[
\int_ {- \infty} ^ {+ \infty} \mathbf {f} (x) d x
\]

defined in FRV, II, §2, No.~1; this follows from formula (1), and from FRV, I, §1, No.~2, Cor. of Prop. 2.

\subsection{Continuity properties of the integral}

PROPOSITION 8. --- Suppose that E is Hausdorff; let \(\mu\) be a measure on X. The mapping \(\mathbf{f} \mapsto \int \mathbf{f} d\mu\) of \(\mathcal{K}(\mathrm{X};\mathbf{E})\) into \(\widehat{\mathbf{E}}\) (No.~3, Cor. 1 of Prop. 7) is the unique continuous linear mapping \(\Phi\) such that \(\Phi(g \cdot \mathbf{a}) = \mu(g) \cdot \mathbf{a}\) for every vector \(\mathbf{a} \in \mathbf{E}\) and every function \(g \in \mathcal{K}(\mathrm{X};\mathbf{C})\) .

To prove the continuity of the mapping \(f \mapsto \int f d\mu\) , it suffices to show that its restriction to \(\mathcal{K}(X, K; E)\) is continuous for every compact subset K of X (TVS, II, §4, No.~4, Prop. 5). We note that if the topology of E is defined by a family of semi-norms \((q_{\alpha})\) , that of \(\mathcal{K}(X, K; E)\) is defined by the family of semi-norms

\[
p _ {\alpha} (\mathbf {f}) = \sup _ {x \in K} q _ {\alpha} \bigl (\mathbf {f} (x) \bigr).
\]

Now, let h be a continuous mapping of X into \([0,1]\) , with compact support and such that \(h(x)=1\) on K; by No.~2, Prop. 6 we have, for every function \(\mathbf{f}\in\mathcal{K}(\mathrm{X},\mathrm{K};\mathrm{E})\) ,

\[
q _ {\alpha} \left(\int \mathbf {f} d \mu\right) = q _ {\alpha} \left(\int h \mathbf {f} d \mu\right) \leqslant \int h (x) q _ {\alpha} (\mathbf {f} (x)) d | \mu | (x) \leqslant | \mu | (h) \cdot p _ {\alpha} (\mathbf {f})
\]

(the \(q_{\alpha}\) being extended by continuity to \(\widehat{\mathbf{E}}\) ), which proves the continuity of \(\mathbf{f} \mapsto \int \mathbf{f} d\mu\) . On the other hand, with the notations of the statement,

\[
\int \left(g (x) \cdot {\bf a}\right) d \mu (x) = \mu (g) \cdot {\bf a}
\]

by virtue of No.~1, Example 1 and Prop. 2 of No.~2 applied to the canonical injection \(\mathbf{C} \cdot \mathbf{a} \to \mathrm{E}\) . Moreover, the subspace of \(\mathcal{K}(\mathrm{X};\mathrm{E})\) formed by the linear combinations \(\sum_{i} g_i \cdot \mathbf{a}_i\) , where \(\mathbf{a}_i \in \mathrm{E}\) and \(g_i \in \mathcal{K}(\mathrm{X};\mathbf{C})\) , is dense in \(\mathcal{K}(\mathrm{X};\mathrm{E})\) (§1, No.~2, Prop. 5), which completes the proof.

PROPOSITION 9. --- Suppose that E is Hausdorff; let f be a continuous mapping of X into E with compact support. When the space \(\mathcal{M}(\mathrm{X};\mathbf{C})\) is equipped with the topology of strictly compact convergence (§1, No.~10), the mapping \(\mu \mapsto \int \mathbf{f} d\mu\) of \(\mathcal{M}(\mathrm{X};\mathbf{C})\) into \(\widehat{\mathrm{E}}\) is the unique continuous linear mapping \(\Psi\) such that \(\Psi(\varepsilon_x) = \mathbf{f}(x)\) for all \(x \in \mathrm{X}\) .

For every \(\mathbf{z}' \in \mathrm{E}'\) ,

\[
\left\langle \int \mathbf {f} d \varepsilon_ {x}, \mathbf {z} ^ {\prime} \right\rangle = \int (\mathbf {z} ^ {\prime} \circ \mathbf {f}) d \varepsilon_ {x} = \mathbf {z} ^ {\prime} (\mathbf {f} (x)) = \langle \mathbf {f} (x), \mathbf {z} ^ {\prime} \rangle ,
\]

whence \(\int \mathbf{f} d\varepsilon_x = \mathbf{f}(x)\) . We know, moreover, that the set of point measures is total in \(\mathcal{M}(\mathrm{X};\mathbf{C})\) for the topology of strictly compact convergence (§2, No.~4, Cor. 4 of Th. 1). Thus it all comes down to proving the continuity of the linear mapping \(u: \mu \mapsto \int \mathbf{f} d\mu\) . For this, consider the linear mapping \(v: \mathbf{z}' \mapsto \langle \mathbf{f}, \mathbf{z}' \rangle\) of \(\mathrm{E}'\) into \(\mathcal{K}(\mathrm{X};\mathbf{C})\) , and let us show that the image under \(v\) of an equicontinuous subset H of \(\mathrm{E}'\) is contained in a strictly compact subset of \(\mathcal{K}(\mathrm{X};\mathbf{C})\) . For, if K is the support of \(\mathbf{f}\) , the functions \(\langle \mathbf{f}, \mathbf{z}' \rangle\) for \(\mathbf{z}' \in \mathrm{H}\) have support contained in K; on the other hand, these functions form an equicontinuous set, and for each \(x \in \mathrm{X}\) the set of \(\mathbf{z}'(\mathbf{f}(x))\) is bounded; our assertion therefore follows from Ascoli's theorem (GT, X, §2, No.~5, Cor. 3 of Th. 2). Now, it follows from formula (1) of No.~1 that \(u\) is none other than the restriction to \(\mathcal{M}(\mathrm{X};\mathbf{C})\) of the transpose \(^t v\) (in the algebraic sense); its continuity therefore follows from the foregoing (TVS, IV, §1, No.~3, Prop. 6).

COROLLARY. --- With hypotheses and notations as in Prop. 9, the restriction of the mapping \(\mu \mapsto \int \mathbf{f} d\mu\) to the set \(\mathcal{M}_{+}(\mathrm{X})\) of positive measures, or to a vaguely bounded subset B of \(\mathcal{M}(\mathrm{X};\mathbf{C})\) , is vaguely continuous.

For, it follows from §1, No.~10, Props. 17 and 18 that, on \(\mathcal{M}_{+}(X)\) or on B, the topology induced by the topology of strictly compact convergence is the same as the topology induced by the vague topology.

\pandocbounded{\includegraphics[keepaspectratio]{images/part001_df9b2198.jpg}}

However, the mapping \(\mu \mapsto \int \mathbf{f}d\mu\) is not necessarily continuous on all of \(\mathcal{M}(\mathrm{X};\mathbf{C})\) for the vague topology (Exer. 2).

\section{PRODUCTS OF MEASURES}

\subsection{The product of two measures}

THEOREM 1. --- Let X and Y be two locally compact spaces, \(\lambda\) a measure on X, \(\mu\) a measure on Y; there exists one and only one measure \(\nu\) on X \(\times\) Y such that, for every function \(g \in \mathcal{K}(\mathrm{X};\mathbf{C})\) and every function \(h \in \mathcal{K}(\mathrm{Y};\mathbf{C})\) ,

\[
\int g (x) h (y) d \nu (x, y) = \left(\int g (x) d \lambda (x)\right) \left(\int h (y) d \mu (y)\right).\tag{1}
\]

Lemma 1. --- Let X and Y be two locally compact spaces, K (resp. L) a compact subset of X (resp. Y).

\begin{enumerate}
\def\labelenumi{(\roman{enumi})}
\tightlist
\item
  The restriction to \(\mathcal{K}(\mathrm{X}\times \mathrm{Y},\mathrm{K}\times \mathrm{L};\mathbf{C})\) of the canonical bijection
\end{enumerate}

\[
\omega : \mathcal {F} (\mathrm{X} \times \mathrm{Y}; \mathbf {C}) \rightarrow \mathcal {F} \bigl (\mathrm{X}; \mathcal {F} (\mathrm{Y}; \mathbf {C}) \bigr)
\]

(S, R, §4, No.~14) is an isometry of the Banach space \(\mathcal{K}(\mathrm{X}\times\mathrm{Y},\mathrm{K}\times\mathrm{L};\mathbf{C})\) onto the Banach space \(\mathcal{K}(\mathrm{X},\mathrm{K};\mathcal{K}(\mathrm{Y},\mathrm{L};\mathbf{C}))\) .

\begin{enumerate}
\def\labelenumi{(\roman{enumi})}
\setcounter{enumi}{1}
\tightlist
\item
  The vector space \(\mathcal{K}(\mathrm{X},\mathrm{K};\mathbf{C})\otimes_{\mathbf{C}}\mathcal{K}(\mathrm{Y},\mathrm{L};\mathbf{C})\) , identified canonically with a subspace of \(\mathcal{K}(\mathrm{X}\times \mathrm{Y},\mathrm{K}\times \mathrm{L};\mathbf{C})\) (A, II, §7, No.~7, comments following the Cor. of Prop. 15), is dense in this Banach space.
\end{enumerate}

It is immediate that the image under \(\omega\) of \(\mathcal{K}(X\times Y,K\times L;C)\) is contained in \(\mathcal{K}(X,K;\mathcal{K}(Y,L;C))\) . Conversely, if u is a continuous mapping of X into \(\mathcal{K}(Y,L;C)\) , with support contained in K, then the mapping \((x,y)\mapsto\mathbf{u}(x)(y)\) of \(X\times Y\) into C is continuous and has support contained in \(K\times L\) , therefore the restriction of \(\omega\) to \(\mathcal{K}(X\times Y,K\times L;C)\) is a bijection of this space onto \(\mathcal{K}(X,K;\mathcal{K}(Y,L;C))\) ; the fact that this restriction is a Banach space isometry follows from the relation

\[
\sup _ {(x, y) \in \mathrm{K} \times \mathrm{L}} | f (x, y) | = \sup _ {x \in \mathrm{K}} \left(\sup _ {y \in \mathrm{L}} | f (x, y) |\right).
\]

This proves (i); on the other hand the image under \(\omega\) , of

\[
\mathcal {K} (\mathrm{X}, \mathrm{K}; \mathbf {C}) \otimes_ {\mathbf {C}} \mathcal {K} (\mathrm{Y}, \mathrm{L}; \mathbf {C})
\]

identified with a subspace of \(\mathcal{K}(\mathrm{X}\times \mathrm{Y},\mathrm{K}\times \mathrm{L};\mathbf{C})\) , is again the space \(\mathcal{K}(\mathrm{X},\mathrm{K};\mathbf{C})\otimes_{\mathbf{C}}\mathcal{K}(\mathrm{Y},\mathrm{L};\mathbf{C})\) but this time identified canonically with a space of mappings of X into \(\mathcal{K}(\mathrm{Y},\mathrm{L};\mathbf{C})\) (A, II, §7, No.~7, Cor. of Prop. 15); but this subspace of \(\mathcal{K}(\mathrm{X},\mathrm{K};\mathcal{K}(\mathrm{Y},\mathrm{L};\mathbf{C}))\) is known to be dense in the latter space (§1, No.~2, Prop. 5), thus the conclusion of (ii) follows from the fact that the restriction of \(\omega\) is a topological isomorphism.

Having proved the lemma, we now observe that every compact subset of \(X \times Y\) is contained in a product \(K \times L\) , where K (resp. L) is a compact subset of X (resp. Y). It therefore follows from Lemma 1, (ii) that the subspace \(\mathcal{K}(X; \mathbf{C}) \otimes_{\mathbf{C}} \mathcal{K}(Y; \mathbf{C})\) is dense in \(\mathcal{K}(X \times Y; \mathbf{C})\) ; since the relation (1) may also be written as \(\nu(g \otimes h) = \lambda(g)\mu(h)\) for \(g \in \mathcal{K}(X; \mathbf{C})\) , \(h \in \mathcal{K}(Y; \mathbf{C})\) , the uniqueness of \(\nu\) follows at once. To prove the existence of \(\nu\) , we shall make use of the following lemma:

Lemma 2. --- With notations as in Lemma 1, for every function \(f \in \mathcal{K}(\mathrm{X} \times \mathrm{Y}, \mathrm{K} \times \mathrm{L}; \mathbf{C})\) the function

\[
y \mapsto h (y) = \int f (x, y) d \lambda (x)\tag{2}
\]

belongs to \(\mathcal{K}(\mathrm{Y},\mathrm{L};\mathbf{C})\)

For every function \(\mathbf{u} \in \mathcal{K}(\mathrm{X};\mathcal{K}(\mathrm{Y},\mathrm{L};\mathbf{C}))\) , the integral \(\int \mathbf{u}(x)d\lambda (x)\) belongs to \(\mathcal{K}(\mathrm{Y},\mathrm{L};\mathbf{C})\) since the latter is a Banach space ( \(\S 3\) , No.~3, Cor. 1 of Prop. 7); but for \(\mathbf{u} = \omega(f)\) and for every \(y \in \mathrm{Y}\) ,

\[
\left\langle \int \mathbf {u} (x) d \lambda (x), \varepsilon_ {y} \right\rangle = \int \mathbf {u} (x) (y) d \lambda (x) = \int f (x, y) d \lambda (x),
\]

whence the lemma.

Consider, then, for every function \(f \in \mathcal{K}(\mathrm{X} \times \mathrm{Y}; \mathbf{C})\) , the number \(\nu(f) = \mu\left(\int f(x, y) d\lambda(x)\right)\) (which we shall also write \(\int d\mu(y) \int f(x, y) d\lambda(x)\) by an abuse of notation); if K (resp. L) is a compact subset of X (resp. Y), then there exists a number \(a_{K}\) (resp. \(b_{L}\) ) such that, for every function \(g \in \mathcal{K}(\mathrm{X}, \mathrm{K}; \mathbf{C})\) (resp. \(h \in \mathcal{K}(\mathrm{Y}, \mathrm{L}; \mathbf{C})\) ), we have \(|\lambda(g)| \leqslant a_{\mathrm{K}} \|g\|\) (resp. \(|\mu(h)| \leqslant b_{\mathrm{L}} \|h\|\) ). It follows that for every function \(f \in \mathcal{K}(\mathrm{X} \times \mathrm{Y}, \mathrm{K} \times \mathrm{L}; \mathbf{C})\) ,

\[
\left| \int f (x, y) d \lambda (x) \right| \leqslant a _ {K} \| f \|
\]

for every \(y \in Y\) , whence \(|\nu(f)| \leqslant a_{K}b_{L}\|f\|\) . The linear form \(\nu\) on \(\mathcal{K}(X \times Y; \mathbf{C})\) is thus a measure on \(X \times Y\) that obviously satisfies (1), which completes the proof of Th. 1.

DEFINITION 1. --- Given two measures \(\lambda\) , \(\mu\) defined, respectively, on two locally compact spaces X, Y, the product measure of \(\lambda\) by \(\mu\) is defined to be the unique measure \(\nu\) on \(X \times Y\) satisfying the relation (1) for every function \(g \in \mathcal{K}(X; \mathbf{C})\) and every function \(h \in \mathcal{K}(Y; \mathbf{C})\) .

In the proof of Th. 1, the roles of the spaces X and Y may be interchanged; canonically identifying \(Y \times X\) and \(X \times Y\) , we thus define on \(X \times Y\) a measure

\[
f \mapsto \int d \lambda (x) \int f (x, y) d \mu (y)
\]

that again satisfies condition (1). We have thus proved the following theorem:

THEOREM 2. --- Let \(\lambda\) , \(\mu\) be two measures defined, respectively, on two locally compact spaces \(X\) , \(Y\) . For every function \(f\) in \(\mathcal{K}(X \times Y; \mathbf{C})\) , the integral of \(f\) with respect to the product measure \(\nu\) of \(\lambda\) by \(\mu\) has the value

\[
\begin{array}{c} \int f (x, y) d \nu (x, y) = \int d \lambda (x) \int f (x, y) d \mu (y) \\ = \int d \mu (y) \int f (x, y) d \lambda (x). \end{array}\tag{3}
\]

Because of the last formula, the integral of f with respect to the product measure \(\nu\) is usually denoted \(\iint f d\lambda d\mu\) , or \(\iint f d\mu d\lambda\) , or \(\iint f \lambda \mu\) , or \(\iint f \mu \lambda\) , or \(\iint f(x, y) d\lambda(x) d\mu(y)\) , or \(\iint f(x, y) d\mu(y) d\lambda(x)\) , or \(\iint f(x, y) \lambda(x)\mu(y)\) , or \(\iint f(x, y) \mu(y)\lambda(x)\) ; it is said to be the double integral of f with respect to \(\lambda\) and \(\mu\) . With this notation, the formula (3) may be written

\[
\begin{array}{c} \iint f (x, y) d \lambda (x) d \mu (y) = \int d \lambda (x) \int f (x, y) d \mu (y) \\ = \int d \mu (y) \int f (x, y) d \lambda (x). \end{array}\tag{4}
\]

Formula (3) shows that if \(\lambda\) and \(\mu\) are real (resp. positive) measures, then the product measure \(\nu\) is real (resp. positive).

Examples. --- 1) The product measure of the Dirac measures \(\varepsilon_{x}\) ( \(x \in X\) ) and \(\varepsilon_{y}\) ( \(y \in Y\) ) is the Dirac measure \(\varepsilon_{(x,y)}\) .

\begin{enumerate}
\def\labelenumi{\arabic{enumi})}
\setcounter{enumi}{1}
\tightlist
\item
  Let us take \(\mathbf{X} = \mathbf{Y} = \mathbf{R}\) , and for \(\lambda\) and \(\mu\) the Lebesgue measure on \(\mathbf{R}\) (§1, No.~3); their product is called Lebesgue measure on \(\mathbf{R}^2\) ; the integral of a function \(f \in \mathcal{K}(\mathbf{R}^2; \mathbf{C})\) with respect to this measure is denoted \(\iint f(x, y) dx dy\) or \(\iint f(x, y) dy dx\) ; formula (4), for a function that is zero outside a compact rectangle \([a, b] \times [c, d]\) , yields the formula
\end{enumerate}

\[
\int_ {c} ^ {d} d y \int_ {a} ^ {b} f (x, y) d x = \int_ {a} ^ {b} d x \int_ {c} ^ {d} f (x, y) d y
\]

proved in FRV, II, §3, No.~6.

Since Lebesgue measure on \(\mathbf{R}\) is invariant under every translation (§1, No.~3), it follows at once that Lebesgue measure on \(\mathbf{R}^2\) is invariant under every translation of \(\mathbf{R}^2\) .

Remark. --- Let E be a Hausdorff locally convex space, and let f be a mapping in \(\mathcal{K}(X \times Y; E)\) such that \(\mathbf{f}(X \times Y)\) is contained in a complete convex subset C of E. Then, for every \(y \in Y\) , the integral \(\mathbf{h}(y) = \int \mathbf{f}(x, y) d\lambda(x)\) belongs to E (§3, No.~3, Prop. 7); moreover, the function h belongs to \(\widetilde{\mathcal{K}}(Y; E)\) : indeed, for every \(z' \in E'\) we have

\[
\langle \mathbf {h} (y), \mathbf {z} ^ {\prime} \rangle = \int \langle \mathbf {f} (x, y), \mathbf {z} ^ {\prime} \rangle d \lambda (x),
\]

therefore \(y \mapsto \langle \mathbf{h}(y), \mathbf{z}' \rangle\) belongs to \(\mathcal{K}(\mathrm{Y}; \mathbf{C})\) by Lemma 2. The integral \(\int h d\mu\) is therefore defined (and a priori belongs to \(E'^{*}\) ); let us show that

\[
\begin{array}{r l} \iint \mathbf {f} (x, y) d \lambda (x) d \mu (y) & = \int d \mu (y) \int \mathbf {f} (x, y) d \lambda (x) \\ & = \int d \lambda (x) \int \mathbf {f} (x, y) d \mu (y), \end{array}\tag{5}
\]

thus generalizing the formula (4). Indeed, for every \(\mathbf{z}' \in \mathrm{E}'\) we have

\[
\begin{array}{r l} \left\langle \iint \mathbf {f} d \lambda d \mu , \mathbf {z} ^ {\prime} \right\rangle & = \iint \langle \mathbf {f}, \mathbf {z} ^ {\prime} \rangle d \lambda d \mu = \int d \mu \int \langle \mathbf {f}, \mathbf {z} ^ {\prime} \rangle d \lambda \\ & = \int \left\langle \int \mathbf {f} d \lambda , \mathbf {z} ^ {\prime} \right\rangle d \mu = \left\langle \int d \mu \int \mathbf {f} d \lambda , \mathbf {z} ^ {\prime} \right\rangle \end{array}
\]

by (4), whence (5).

\subsection{Properties of product measures}

If \(\lambda\) (resp. \(\mu\) ) is a measure on X (resp. Y) and \(\nu\) is the product measure of \(\lambda\) by \(\mu\) , then the restriction of \(\nu\) to \(\mathcal{K}(X;C)\otimes_{\mathbf{C}}\mathcal{K}(Y;C)\) is none other than the tensor product \(\lambda\otimes\mu\) of the two linear forms \(\lambda\) and \(\mu\) (A, II, §3, No.~2), because the relation (1) of No.~1 may be written

\[
\langle g \otimes h, \nu \rangle = \langle g, \lambda \rangle \langle h, \mu \rangle = \langle g \otimes h, \lambda \otimes \mu \rangle
\]

for all \(g \in \mathcal{K}(X; \mathbf{C})\) and \(h \in \mathcal{K}(Y; \mathbf{C})\) . By an abuse of notation, we shall also denote by \(\lambda \otimes \mu\) the product measure \(\nu\) (and not just its restriction to the dense subspace \(\mathcal{K}(X; \mathbf{C}) \otimes_{\mathbf{C}} \mathcal{K}(Y; \mathbf{C})\) of \(\mathcal{K}(X \times Y; \mathbf{C})\) ).

The mapping \((\lambda, \mu) \mapsto \lambda \otimes \mu\) of \(\mathcal{M}(\mathrm{X}; \mathbf{C}) \times \mathcal{M}(\mathrm{Y}; \mathbf{C})\) into \(\mathcal{M}(\mathrm{X} \times \mathrm{Y}; \mathbf{C})\) is obviously bilinear, by virtue of formula (3) of No.~1.

PROPOSITION 1. --- Let \(\lambda\) be a measure on X, \(\mu\) a measure on Y; if \(g \in \mathcal{C}(\mathrm{X};\mathbf{C})\) , \(h \in \mathcal{C}(\mathrm{Y};\mathbf{C})\) , then

\[
(g \cdot \lambda) \otimes (h \cdot \mu) = (g \otimes h) \cdot (\lambda \otimes \mu).\tag{6}
\]

For every function \(f \in \mathcal{K}(\mathrm{X} \times \mathrm{Y}; \mathbf{C})\) , we have, by virtue of formula (3) of No.~1,

\[
\begin{array}{c} \langle f, (g \otimes h) \cdot (\lambda \otimes \mu) \rangle = \int d \lambda (x) \int f (x, y) g (x) h (y) d \mu (y) \\ = \int g (x) d \lambda (x) \int f (x, y) h (y) d \mu (y), \end{array}
\]

which proves formula (6).

PROPOSITION 2. --- The support of the product \(\lambda \otimes \mu\) is equal to the product of the support of \(\lambda\) and the support of \(\mu\) .

We first observe that the relation \(\lambda \otimes \mu = 0\) implies that one of the measures \(\lambda, \mu\) is zero (A, II, §7, No.~7, Prop. 16, (ii)). On the other hand, if U (resp. V) is an open set in X (resp. Y), then the restriction of \(\lambda\otimes\mu\) to the product \(U\times V\) is the product of the restrictions of \(\lambda\) to U and of \(\mu\) to V, as follows from Th. 1 of No.~1 and the definition of the restriction of a measure to an open set (§2, No.~1). It therefore follows that, for the restriction of \(\lambda\otimes\mu\) to \(U\times V\) to be zero, it is necessary and sufficient that either the restriction of \(\lambda\) to U or the restriction of \(\mu\) to V be zero, which proves the proposition, on taking into account the definition of the support of a measure (§2, No.~2).

PROPOSITION 3. --- Let \(\lambda \in \mathcal{M}(\mathrm{X};\mathbf{C})\) , \(\mu \in \mathcal{M}(\mathrm{Y};\mathbf{C})\) . Then

\[
| \lambda \otimes \mu | = | \lambda | \otimes | \mu |.\tag{7}
\]

Let \(f \in \mathcal{K}_{+}(\mathrm{X} \times \mathrm{Y})\) , \(g \in \mathcal{K}(\mathrm{X} \times \mathrm{Y}; \mathbf{C})\) be such that \(|g| \leqslant f\) ; we have (§1, No.~6, formula (13))

\[
\begin{array}{r l} | \langle g, \lambda \otimes \mu \rangle | = & \left| \int d \lambda (x) \int g (x, y)   d \mu (y) \right| \\ & \leqslant \int d | \lambda | (x) \int | g (x, y) |   d | \mu | (y) \\ & = \langle | g |, | \lambda | \otimes | \mu | \rangle \leqslant \langle f, | \lambda | \otimes | \mu | \rangle . \end{array}
\]

It follows that \(\langle f, |\lambda \otimes \mu| \rangle \leqslant \langle f, |\lambda| \otimes |\mu| \rangle\) , and so

\[
| \lambda \otimes \mu | \leqslant | \lambda | \otimes | \mu |.\tag{8}
\]

On the other hand, let \(u \in \mathcal{K}_{+}(\mathrm{X})\) , \(v \in \mathcal{K}_{+}(\mathrm{Y})\) . For every \(\varepsilon > 0\) , there exist \(u_{1} \in \mathcal{K}(\mathrm{X};\mathbf{C})\) , \(v_{1} \in \mathcal{K}(\mathrm{Y};\mathbf{C})\) such that \(|u_{1}| \leqslant u\) , \(|v_{1}| \leqslant v\) and

\[
| \langle u _ {1}, \lambda \rangle | \geqslant \langle u, | \lambda | \rangle - \varepsilon , \quad | \langle v _ {1}, \mu \rangle | \geqslant \langle v, | \mu | \rangle - \varepsilon
\]

(§1, No.~6). It follows that \(|u_1 \otimes v_1| \leqslant u \otimes v\) and

\[
\begin{array}{c} \langle u \otimes v, | \lambda \otimes \mu | \rangle \geqslant | \langle u _ {1} \otimes v _ {1}, \lambda \otimes \mu \rangle | = | \langle u _ {1}, \lambda \rangle \langle v _ {1}, \mu \rangle | \\ \geqslant (\langle u, | \lambda | \rangle - \varepsilon) (\langle v, | \mu | \rangle - \varepsilon). \end{array}
\]

Since \(\varepsilon\) is arbitrary, we infer that

\[
\langle u \otimes v, | \lambda \otimes \mu | \rangle \geqslant \langle u, | \lambda | \rangle \langle v, | \mu | \rangle = \langle u \otimes v, | \lambda | \otimes | \mu | \rangle .
\]

Taking (8) into account, we see that

\[
\langle u \otimes v, | \lambda \otimes \mu | \rangle = \langle u \otimes v, | \lambda | \otimes | \mu | \rangle .
\]

Since every function in \(\mathcal{K}(\mathrm{X};\mathbf{C})\) (resp. \(\mathcal{K}(\mathrm{Y};\mathbf{C})\) ) is a linear combination of functions in \(\mathcal{K}_{+}(\mathrm{X})\) (resp. \(\mathcal{K}_{+}(\mathrm{Y})\) ), the preceding formula remains true for \(u\in\mathcal{K}(\mathrm{X};\mathbf{C})\) and \(v\in\mathcal{K}(\mathrm{Y};\mathbf{C})\) ; the proposition therefore follows from the fact that \(\mathcal{K}(\mathrm{X};\mathbf{C})\otimes_{\mathbf{C}}\mathcal{K}(\mathrm{Y};\mathbf{C})\) is dense in \(\mathcal{K}(\mathrm{X}\times\mathrm{Y};\mathbf{C})\) .

COROLLARY. --- Let \(\lambda \in \mathcal{M}(\mathrm{X};\mathbf{R})\) , \(\mu \in \mathcal{M}(\mathrm{Y};\mathbf{R})\) . Then

\[
\left\{ \begin{array}{l} (\lambda \otimes \mu) ^ {+} = \lambda^ {+} \otimes \mu^ {+} + \lambda^ {-} \otimes \mu^ {-}, \\ (\lambda \otimes \mu) ^ {-} = \lambda^ {+} \otimes \mu^ {-} + \lambda^ {-} \otimes \mu^ {+}. \end{array} \right.\tag{9}
\]

For, by virtue of Prop. 3,

\[
\begin{array}{r l} & {(\lambda \otimes \mu) ^ {+} = \frac {1}{2} \big (\lambda \otimes \mu + | \lambda | \otimes | \mu | \big)} \\ & {\qquad = \frac {1}{2} \big ((\lambda^ {+} - \lambda^ {-}) \otimes (\mu^ {+} - \mu^ {-}) + (\lambda^ {+} + \lambda^ {-}) \otimes (\mu^ {+} + \mu^ {-}) \big)} \\ & {\qquad = \lambda^ {+} \otimes \mu^ {+} + \lambda^ {-} \otimes \mu^ {-}.} \end{array}
\]

The argument for \((\lambda \otimes \mu)^{-}\) is similar.

PROPOSITION 4. --- Let \(\lambda \in \mathcal{M}(\mathrm{X};\mathbf{C})\) , \(\mu \in \mathcal{M}(\mathrm{Y};\mathbf{C})\) . Then

\[
\| \lambda \otimes \mu \| = \| \lambda \| \cdot \| \mu \|,\tag{10}
\]

with the convention that the second member is to be replaced by 0 whenever one of the factors is 0 and the other is +∞. In particular, if λ and μ are bounded then λ⊗μ is bounded.

By the above Proposition 3, and §1, No.~8, Cor. 1 of Prop. 10, we may limit ourselves to the case that \(\lambda\) and \(\mu\) are positive measures. If \(\lambda = 0\) or \(\mu = 0\) , the result is trivial; let us therefore assume that \(\lambda \neq 0\) and \(\mu \neq 0\) . Let us first prove that \(\| \lambda \otimes \mu \| \leqslant \| \lambda \| \cdot \| \mu \|\) . We may suppose \(\lambda\) and \(\mu\) to be bounded. For every \(f \in \mathcal{K}_+(\mathrm{X} \times \mathrm{Y})\) ,

\[
\langle f, \lambda \otimes \mu \rangle = \int d \lambda (x) \int f (x, y) d \mu (y)
\]

and

\[
\int f (x, y) d \mu (y) \leqslant \| f \| \cdot \| \mu \|
\]

for every \(x\in \mathbf{X}\) , therefore

\[
\langle f, \lambda \otimes \mu \rangle \leqslant \| f \| \cdot \| \lambda \| \cdot \| \mu \|,
\]

which proves our assertion. On the other hand, let

\[
\alpha <   \| \lambda \|, \beta <   \| \mu \|
\]

be two real numbers \(\geqslant 0\) . There exist \(g\in \mathcal{K}_{+}(\mathrm{X})\) , \(h\in \mathcal{K}_{+}(\mathrm{Y})\) such that

\[
\| g \| \leqslant 1, \| h \| \leqslant 1, \lambda (g) \geqslant \alpha , \mu (h) \geqslant \beta .
\]

Then \(g \otimes h \in \mathcal{K}_{+}(X \times Y)\) , \(\| g \otimes h \| \leqslant 1\) and \(\langle g \otimes h, \lambda \otimes \mu \rangle \geqslant \alpha \beta\) ; therefore \(\| \lambda \otimes \mu \| \geqslant \alpha \beta\) and finally \(\| \lambda \otimes \mu \| \geqslant \| \lambda \| \cdot \| \mu \|\) , which completes the proof.

\subsection{Continuity of product measures}

PROPOSITION 5. --- For every measure \(\lambda_0 \in \mathcal{M}(\mathrm{X};\mathbf{C})\) , the mapping \(\mu \mapsto \lambda_0 \otimes \mu\) of \(\mathcal{M}(\mathrm{Y};\mathbf{C})\) into \(\mathcal{M}(\mathrm{X} \times \mathrm{Y};\mathbf{C})\) is vaguely continuous.

For every function \(f \in \mathcal{K}(\mathrm{X} \times \mathrm{Y}; \mathbf{C})\) , we know that the function \(h(y) = \int f(x, y) d\lambda_0(x)\) belongs to \(\mathcal{K}(\mathrm{Y}; \mathbf{C})\) (No.~1, Lemma 2), and \(\langle f, \lambda_0 \otimes \mu \rangle = \langle h, \mu \rangle\) , whence the proposition.

PROPOSITION 6. --- When \(\mathcal{M}(\mathrm{X};\mathbf{C})\) , \(\mathcal{M}(\mathrm{Y};\mathbf{C})\) and \(\mathcal{M}(\mathrm{X}\times\mathrm{Y};\mathbf{C})\) are equipped with the topology of strictly compact convergence (§1, No.~10), the bilinear mapping \((\lambda,\mu)\mapsto\lambda\otimes\mu\) of \(\mathcal{M}(\mathrm{X};\mathbf{C})\times\mathcal{M}(\mathrm{Y};\mathbf{C})\) into \(\mathcal{M}(\mathrm{X}\times\mathrm{Y};\mathbf{C})\) is hypocontinuous for the set of vaguely bounded subsets of \(\mathcal{M}(\mathrm{X};\mathbf{C})\) and \(\mathcal{M}(\mathrm{Y};\mathbf{C})\) (TVS, III, §5, No.~3).

Let \(K \subset X\) , \(L \subset Y\) be two compact sets, A a compact subset of \(\mathcal{K}(X \times Y, K \times L; C)\) , and B a vaguely bounded and closed subset of \(\mathcal{M}(X; C)\) ; it is known that B is vaguely compact (§1, No.~9, Prop. 15), hence also compact for the topology of strictly compact convergence (§1, No.~10, Prop. 17). On the other hand, the Banach space \(\mathcal{K}(X \times Y, K \times L; C)\) is isometric to \(\mathcal{K}(X, K; \mathcal{K}(Y, L; C))\) (No.~1, Lemma 1); the mapping \(\varphi\) of \(\mathcal{K}(X, K; \mathcal{K}(Y, L; C)) \times \mathcal{M}(X; C)\) into \(\mathcal{K}(Y, L; C)\) , such that \(\varphi(g, \lambda)\) is the function h defined by \(h(y) = \int g(x, y) d\lambda(x)\) , is separately continuous by virtue of §3, No.~4, Props. 8 and 9. Since \(\mathcal{K}(X, K; \mathcal{K}(Y, L; C))\) is barreled, it follows that the mapping \(\varphi\) is hypocontinuous relative to the vaguely bounded subsets of \(\mathcal{M}(X; C)\) (TVS, III, §5, No.~3, Prop. 6); the restriction of this mapping to \(A \times B\) is therefore continuous (loc. cit., Prop. 4). The image C of \(A \times B\) under this mapping is consequently a compact subset of the Banach space \(\mathcal{K}(Y, L; C)\) . Now, C is none other than the set of functions \(h(y) = \int f(x, y) d\lambda(x)\) as f runs over A and \(\lambda\) runs over B; by virtue of formula (3) of No.~1, the conditions \(\lambda \in B\) and \(\mu \in C^{\circ}\) therefore imply \(\lambda \otimes \mu \in A^{\circ}\) . In view of the definition of the topology of strictly compact convergence, this proves the proposition (TVS, III, §5, No.~3, Def. 2).

\pandocbounded{\includegraphics[keepaspectratio]{images/part001_f639fdbe.jpg}}

The conclusion of Prop. 6 is no longer valid when the topology of strictly compact convergence is replaced by the vague topology (Exer. 2 c)). However, if B (resp. C) is a vaguely bounded subset of \(\mathcal{M}(\mathrm{X};\mathbf{C})\) (resp. \(\mathcal{M}(\mathrm{Y};\mathbf{C})\) ), then the image of \(\mathrm{B}\times \mathrm{C}\) under the mapping \((\lambda ,\mu)\mapsto \lambda \otimes \mu\) is vaguely bounded in \(\mathcal{M}(\mathrm{X}\times \mathrm{Y};\mathbf{C})\) and therefore the restriction of this mapping to \(\mathrm{B}\times \mathrm{C}\) is vaguely continuous, by virtue of Prop. 6, of §1, No.~10, Prop. 17, and of Prop. 4 of TVS, III, §5, No.~3 (cf.~Exer. 3).

\subsection{Product of a finite number of measures}

Let \(X_{i}\) ( \(1 \leqslant i \leqslant n\) ) be \(n\) locally compact spaces, \(X = \prod_{i=1}^{n} X_{i}\) their product. The set of linear combinations of complex functions of the form

\[
\left(x _ {1}, x _ {2}, \dots , x _ {n}\right) \mapsto f _ {1} \left(x _ {1}\right) f _ {2} \left(x _ {2}\right) \dots f _ {n} \left(x _ {n}\right),
\]

where \(f_{i} \in \mathcal{K}(\mathrm{X}_{i};\mathbf{C})\) ( \(1 \leqslant i \leqslant n\) ), may be identified canonically with the tensor product \(\bigotimes_{i=1}^{n} \mathcal{K}(\mathrm{X}_{i};\mathbf{C})\) , and it follows from Lemma 1 of No.~1, by induction on \(n\) , that this tensor product is dense in \(\mathcal{K}(\mathrm{X};\mathbf{C})\) .

Now let \(\mu_{i}\) be a measure on \(X_{i}\) ( \(1 \leqslant i \leqslant n\) ); there exists on X one and only one measure \(\nu\) such that, for \(f_{i} \in \mathcal{K}(X_{i}; \mathbf{C})\) ( \(1 \leqslant i \leqslant n\) ),

\[
\left\langle f _ {1} \otimes f _ {2} \otimes \dots \otimes f _ {n}, \nu \right\rangle = \prod_ {i = 1} ^ {n} \left\langle f _ {i}, \mu_ {i} \right\rangle .\tag{11}
\]

For, if this measure exists, it is unique by the foregoing. On the other hand, let \(\nu = \mu_{1} \otimes \mu_{2} \otimes \cdots \otimes \mu_{n}\) be the measure on X defined by the recursion relation

\[
\mu_ {1} \otimes \mu_ {2} \otimes \dots \otimes \mu_ {n} = (\mu_ {1} \otimes \mu_ {2} \otimes \dots \otimes \mu_ {n - 1}) \otimes \mu_ {n}.
\]

It follows from No.~1, formula (1) and this definition (by induction on n) that \(\nu\) verifies (11); it is said to be the product measure of the measures \(\mu_{i}\) ( \(1 \leqslant i \leqslant n\) ) and it is denoted again by \(\bigotimes_{i=1}^{n} \mu_{i}\) . The relation (11) may also be written

\[
\left\langle f _ {1} \otimes f _ {2} \otimes \dots \otimes f _ {n}, \mu_ {1} \otimes \mu_ {2} \otimes \dots \otimes \mu_ {n} \right\rangle = \prod_ {i = 1} ^ {n} \left\langle f _ {i}, \mu_ {i} \right\rangle .\tag{12}
\]

PROPOSITION 7 (`associativity of the product of measures'). --- Let \((\mathbf{I}_k)_{1 \leqslant k \leqslant r}\) be a partition of the interval \([1, n]\) of \(\mathbf{N}\) ; then

\[
\bigotimes_ {k = 1} ^ {r} \left(\bigotimes_ {i \in I _ {k}} \mu_ {i}\right) = \bigotimes_ {i = 1} ^ {n} \mu_ {i}\tag{13}
\]

For, these two measures coincide, by (12), for every function in \(\bigotimes_ {i = 1} ^ {n} \mathcal{K}(X_{i}; \mathbf{C})\) .

The integral of a function \(f \in \mathcal{K}(\mathbf{X};\mathbf{C})\) with respect to the product measure is denoted

\[
\int f d \mu_ {1} d \mu_ {2} \dots d \mu_ {n},
\]

or

\[
\iint \ldots \int f d \mu_ {1} d \mu_ {2} \ldots d \mu_ {n}
\]

or

\[
\int f (\mu_ {1} \otimes \dots \otimes \mu_ {n})
\]

or also

\[
\iint \dots \int f (x _ {1}, x _ {2}, \dots , x _ {n}) d \mu_ {1} (x _ {1}) d \mu_ {2} (x _ {2}) \dots d \mu_ {n} (x _ {n})
\]

or

\[
\iint \dots \int f (x _ {1}, x _ {2}, \dots , x _ {n}) \mu_ {1} (x _ {1}) \mu_ {2} (x _ {2}) \dots \mu_ {n} (x _ {n})
\]

with n signs \(\int\) ; it is said to be a multiple integral of order n, or an n-tuple integral. By virtue of the associativity of the product of measures and the theorem on inverting the order of integration (No.~1, Th. 2), we have, for every permutation \(\sigma\) of \([1,n]\) ,

\[
\iint \dots \int f d \mu_ {1} d \mu_ {2} \dots d \mu_ {n} = \int d \mu_ {\sigma (1)} \int d \mu_ {\sigma (2)} \dots \int f d \mu_ {\sigma (n)}.\tag{14}
\]

The integral notation and formula (14) may be extended in an obvious way to functions \(\mathbf{f} \in \mathcal{K}(\mathrm{X};\mathrm{E})\) with values in a Hausdorff locally convex space E, such that \(\mathbf{f}(\mathrm{X})\) is contained in a complete convex subset of E. We leave to the reader the task of generalizing to the product of any finite number of measures the results of Nos. 2 and 3 concerning the product of two measures.

In particular, one calls Lebesgue measure on \(\mathbf{R}^n\) the product of \(n\) measures identical to the Lebesgue measure on \(\mathbf{R}\) ; the integral of a function \(\mathbf{f} \in \mathcal{K}(\mathbf{R}^n; \mathbf{E})\) , satisfying the preceding condition, is denoted

\[
\iint \dots \int \mathbf {f} (x _ {1}, x _ {2}, \dots , x _ {n}) d x _ {1} d x _ {2} \dots d x _ {n}
\]

and is equal to

\[
\int_ {- \infty} ^ {+ \infty} d x _ {1} \int_ {- \infty} ^ {+ \infty} d x _ {2} \dots \int_ {- \infty} ^ {+ \infty} \mathbf {f} (x _ {1}, x _ {2}, \dots , x _ {n}) d x _ {n}.
\]

Lebesgue measure on \(R^{n}\) is invariant under every translation.
